% =====================================================================
%  Tones Hidden in Tokens: How LLMs Handle Vietnamese Sub-Syllabic Structure
%  ACL 2027 long paper (ARR January 2027 cycle) -- review version
%
%  Build (pdflatex is not installed on the build machine; run on a TeX Live
%  2024+ installation with the vntex package for the T5 encoding):
%      cd paper && pdflatex main && bibtex main && pdflatex main && pdflatex main
%  Every table and figure under tables/ and figures/ is GENERATED by a script in
%  this directory (gen_table1.py, gen_rule_tables.py, gen_fig1_tokens.py,
%  gen_prompt_appendix.py); no number in the paper is typed by hand.
%
%  Page budget (docs/PLAN_2026-09-30.md, section 2.7; 8 content pages):
%      1 Introduction .................. 1.00 p  (Fig. 1)
%      2 Background .................... 0.75 p
%      3 The NoiLai benchmark .......... 1.50 p  (Table 1)
%      4 Setup ......................... 0.75 p
%      5 Results ....................... 1.50 p  (Table 2, Figs. 2-3)
%      6 Orthographic counterfactuals .. 1.00 p  (Table 3, Fig. 4)
%      7 Where tone lives .............. 1.00 p  (Fig. 5)
%      8 Discussion and conclusion ..... 0.50 p
%      Limitations, Ethical considerations, references, appendices: unlimited
% =====================================================================
\documentclass[11pt]{article}

% "review" = anonymous, line numbers, page numbers. Change to "final" for camera-ready.
\usepackage[review]{acl}

\usepackage{times}
\usepackage{latexsym}

% Vietnamese needs the T5 font encoding (vntex, in TeX Live). T5 is a superset of the
% T1 repertoire for Latin text, so it is listed last and becomes the default; T1 stays
% available for the odd glyph. `times` maps rm/sf/tt to ptm/phv/pcr, all of which vntex
% provides in T5 (t5ptm.fd, t5phv.fd, t5pcr.fd). inconsolata (zi4) has no T5 variant and
% is therefore NOT loaded. T5 also defines the accent macros \h (hook above) and \horn,
% which figures/fig1_tokens.tex uses to draw isolated combining marks.
\usepackage[T1,T5]{fontenc}
\usepackage[utf8]{inputenc}

\usepackage{microtype}
\usepackage{graphicx}
\usepackage{booktabs}
\usepackage{multirow}
\usepackage{amsmath}
\usepackage{amssymb}
\usepackage{enumitem}
% xcolor and hyperref are loaded by acl.sty.

% ---------------------------------------------------------------------
% \placeholder{...}: every number or claim that awaits a real result. Red in review
% mode so that nothing pending can be mistaken for a finding; bold brackets in the final
% copy so that a leftover placeholder is still visible and never ships silently.
% tests/test_paper.py checks that the string TODO occurs only inside \placeholder{...}.
% ---------------------------------------------------------------------
\makeatletter
\ifacl@finalcopy
  \newcommand{\placeholder}[1]{\textbf{[#1]}}
\else
  \newcommand{\placeholder}[1]{\textcolor{red}{[#1]}}
\fi
\makeatother

% Hypothesis labels (pre-registered in docs/PREREGISTRATION.md).
\newcommand{\hyp}[1]{\textbf{H#1}}
% Vietnamese examples in running text.
\newcommand{\vi}[1]{\textit{#1}}
% Token boxes for Figure 1 (filled by gen_fig1_tokens.py).
\newcommand{\tok}[1]{\fbox{\strut\small\texttt{#1}}\hspace{1pt}}
\providecommand{\tokspace}{\textvisiblespace}
% Task and variant names.
\newcommand{\task}[1]{\textsc{#1}}
\newcommand{\var}[1]{V#1}
\newcommand{\noilai}{NóiLái}

\title{Tones Hidden in Tokens: How LLMs Handle Vietnamese Sub-Syllabic Structure}

% Anonymized for review. The final version replaces this block with the author list;
% the review option hides it anyway, but nothing identifying is written here on purpose.
\author{Anonymous ACL submission}

\begin{document}
\maketitle

\begin{abstract}
Written Vietnamese encodes each syllable's onset, rime and tone, yet subword tokenizers split syllables without regard to this structure. Can large language models manipulate units they never receive as tokens? We introduce \noilai{}, a benchmark built on \vi{nói lái}, a native Vietnamese wordplay that swaps rimes, onsets or tones between syllables. A rule engine encoding Vietnamese spelling generates \placeholder{$\approx$10,000} automatically verifiable items across three tasks (transformation, decoding, validity judgment) and the four generated variants of a six-type tradition, plus \placeholder{200--400} attested folk and literary examples; native speakers validated a probability sample (\placeholder{$\alpha$ = TODO}). Across \placeholder{17} self-hosted and \placeholder{4} API-served models, \placeholder{\hyp{1}: accuracy tracks how well token boundaries align with onset and rime boundaries more than raw token count; best model TODO vs.\ humans TODO}. We then intervene on tokenization with meaning and information held fixed, re-encoding text between precomposed and combining diacritics and between tone-mark placement conventions, and estimate the within-item causal effects on \noilai{} and on an existing Vietnamese benchmark, testing whether they track the number of extra tokens or the presence of an explicit tone symbol (\placeholder{\hyp{3}: TODO}). Probes and activation patching in Gemma~3 1B and 4B locate \placeholder{\hyp{5}: where tone becomes linearly decodable and whether it is used, TODO}. We release the generator, syllable parser and re-encoding toolkit.
\end{abstract}

% Section 1 -- Introduction. Budget: 1.00 page, holds Figure 1.
\section{Introduction}
\label{sec:intro}

Written Vietnamese is a syllabic script in a Latin alphabet: every syllable is written as a separate word, and its three parts are visible on the page. The onset is a consonant letter or digraph, the rime is a vowel sequence with an optional final consonant, and the tone is one of six values marked by a diacritic (or by its absence). A literate speaker can take a syllable apart at these seams and put it back together, and Vietnamese has a folk art that does exactly that. \vi{Nói lái} swaps the rimes, the tones, or both between two syllables of a phrase: \vi{mèo cái} `female cat' becomes \vi{mài kéo} `sharpen scissors', and \vi{bí mật} `secret' becomes \vi{bị mất} `got lost'. The transformation is defined entirely below the level of the word, and the spelling of the result follows rules that depend on the new neighbours: when the rime \vi{eo} moves next to the onset /k/, the letter changes from \vi{c} to \vi{k}.

Large language models do not see any of this structure. Subword tokenizers split Vietnamese syllables where their merge statistics tell them to, not at the onset--rime boundary, and when text arrives in a different but equally valid Unicode encoding the split changes altogether. Figure~\ref{fig:tokens} shows the Gemma~3 tokenizer on the example above: in the precomposed (NFC) encoding the output syllable \vi{mài} is cut between onset and rime, and in the combining (NFD) encoding the tone mark of \vi{mèo} becomes a token of its own. Our audit of the same tokenizer over the 6,595 standard Vietnamese syllables finds that it splits 70\% of them under NFC, half of all syllables exactly at the onset--rime boundary, and spends one extra token per syllable under NFD, where the tone mark is isolated in 42\% of syllables.\footnote{Numbers from the tokenizer audit in our released code (\texttt{data/audit/gemma3.json}); the tokenizer does not normalize NFD to NFC, so the two encodings reach the model as different token sequences.} The question this paper asks is therefore general: \emph{can a subword LLM reach the onset, rime and tone inside a tonal syllable it never receives as a unit, and does failure grow with how badly the tokens cut across those units?}

Nói lái is our instrument, not our headline. Because the documented variants are permutations of syllable components and Vietnamese spelling is rule-governed, a small rule engine can generate transformation, decoding and validity items in unlimited quantity, verify every gold answer, and score every model output by parsing it back into onset, rime and tone. No LLM judge is needed anywhere. The same engine lets us re-encode any Vietnamese text between canonically equivalent Unicode forms and between the two tone-mark placement conventions while holding meaning and information fixed, which turns deployed tokenizers into the object of a training-free intervention. Finally, the two smallest Gemma~3 models let us ask where in the network tone information is represented, whether it reaches the answer position, and whether it is used, following the template of character-level probing work \citep{2025-findings-emnlp-719}.

% The caption's data sentences (\figtokensnote) are GENERATED by gen_fig1_tokens.py from the
% same audit rows as the table, so that the caption cannot contradict the numbers it sits under.
% GENERATED by paper/gen_fig1_tokens.py from /home/user/noilai-migrate/data/external/gemma3_tokenizer.model on 2026-10-01 13:09Z; do not edit by hand.
% \figtokensnote: the caption sentences of Figure 1 that state what the token table shows,
% composed from the audit rows so that the caption cannot contradict the figure.
\newcommand{\figtokensnote}{Under NFC \vi{mèo}, \vi{cái} and \vi{kéo} are single tokens and \vi{mài} is cut at the onset$|$rime seam (alignment 1.00). Under NFD every syllable gains a boundary between a base letter and its combining tone mark, which is never linguistic, so alignment falls to 0.50 for \vi{mèo} and \vi{kéo}, whose tone mark becomes a token of its own, leaving one aligned boundary at the nucleus$|$coda seam, and to 0.00 for \vi{cái} and \vi{mài}, whose tone mark stays attached to the following letter.}

\begin{figure}[t]
  \centering
  \small
  % GENERATED by paper/gen_fig1_tokens.py from /home/user/ntcf/noilai/data/external/gemma3_tokenizer.model on 2026-09-30 15:10Z; do not edit by hand.
% tokenizer: BPE 262,144 pieces, normalizer identity, byte_fallback True; syllables tokenized in running-text position (leading space).
\begin{tabular}{@{}lllcll@{}}
\toprule
 & \vi{mèo} & \vi{cái} &  & \vi{mài} & \vi{kéo} \\
 & \multicolumn{2}{l}{input (\vi{mèo cái})} & $\xrightarrow{\ \var{1}\ }$ & \multicolumn{2}{l}{output (\vi{mài kéo})} \\
\midrule
NFC (precomposed) & \tok{\tokspace{}mèo} & \tok{\tokspace{}cái} &  & \tok{\tokspace{}m}\tok{ài} & \tok{\tokspace{}kéo} \\
\quad tokens / alignment & 1 / 1.00 & 1 / 1.00 &  & 2 / 1.00 & 1 / 1.00 \\
NFD (combining) & \tok{\tokspace{}me}\tok{\`{}}\tok{o} & \tok{\tokspace{}ca}\tok{\'{}i} &  & \tok{\tokspace{}ma}\tok{\`{}i} & \tok{\tokspace{}ke}\tok{\'{}}\tok{o} \\
\quad tokens / alignment & 3 / 0.50 & 2 / 0.00 &  & 2 / 0.00 & 3 / 0.50 \\
\bottomrule
\end{tabular}

  \caption{\vi{Mèo cái} $\rightarrow$ \vi{mài kéo} (variant \var{1}: the rimes swap, onsets and tones stay) with the token boundaries of the Gemma~3 tokenizer, each syllable in running-text position. \protect\tokspace{} marks the SentencePiece word-initial space; an accent over an empty box is a combining mark that forms a token on its own. Alignment is the share of internal token boundaries that fall on an onset$|$rime, glide$|$nucleus or nucleus$|$coda seam; a boundary between a base letter and its combining mark never counts. \figtokensnote{} \placeholder{Add model outputs for three tokenizers once E1 has run.}}
  \label{fig:tokens}
\end{figure}

We make five contributions.
\begin{enumerate}[nosep,leftmargin=*]
  \item \textbf{\noilai{}}, the first LLM study of nói lái and the first generative, decoding and validity benchmark of it: a rule-verified generator over the four generated variants of a six-type tradition, an attested cultural subset, native-speaker validation and a human baseline (Section~\ref{sec:benchmark}).
  \item \textbf{A syllable-alignment audit} that links how production tokenizers split onset, rime and tone to item-level accuracy across \placeholder{21} models (Sections~\ref{sec:setup} and~\ref{sec:results}).
  \item \textbf{Orthographic counterfactuals} that change only the token sequence of a tonal Latin script on deployed models, separating tokenizer-caused from knowledge-caused errors, with a census of which tokenizers normalize, pass through or corrupt combining marks (Section~\ref{sec:counterfactuals}).
  \item \textbf{A mechanistic localization} of tone information in small open models, using control-validated probes, structural baselines and activation patching with separate perception and manipulation readouts (Section~\ref{sec:tone}).
  \item \textbf{Reusable tooling}: the syllable parser, spelling-rule engine and re-encoder apply to any diacritic-heavy Latin script.
\end{enumerate}

Every headline number carries a 95\% interval that resamples base pairs, every comparison is a paired test corrected within its family, every claim about models in general rests on a pooled estimate rather than a count of significant cells, and the hypotheses \hyp{1}--\hyp{6} named below were pre-registered in a dated commit before any test-split run (the pilot runs on development items only). \placeholder{One-sentence summary of the main finding, written from the results: TODO.}

% Section 2 -- Background: the syllable on the page, three encodings, noi lai. Budget: 0.75 page.
% Follows docs/DESIGN_DECISIONS.md sections 2, 3 (six-type taxonomy, region-neutral labels,
% old-style placement as the majority convention). Related work is Section 8; the full rule
% tables are Appendix B. Claim ledger: paper/claims.md rows L1-L12, N1, D13.
% Every Vietnamese example below is a row of data/attested_seed.tsv or an engine output
% (DESIGN_DECISIONS 3.6); all await native-speaker verification (the NV placeholder).
\section{Background}
\label{sec:background}

\paragraph{The Vietnamese syllable on the page.}
A written syllable is onset + rime + tone, the rime being an optional glide, a nucleus and an optional coda \citep{thompson-1965-vietnamese,doan-1977-nguam,nguyen-1997-vietnamese}; our parser recognizes 24 onsets (including the empty one), 18 nuclei, 11 codas (including none) and six tones, \vi{ngang} (unmarked), \vi{huyền}, \vi{sắc}, \vi{hỏi}, \vi{ngã} and \vi{nặng} \citep{kirby-2011-vietnamese,pham-2003-vietnamese-tone}. Structure and letters are not one-to-one: /k/ is written \vi{q} before the glide, \vi{k} before \vi{e ê i y} and \vi{c} elsewhere, and syllables ending in a stop take only \vi{sắc} and \vi{nặng} (Appendix~\ref{app:rules}). Two conventions place the tone mark; they differ only on the open rimes \vi{oa}, \vi{oe} and \vi{uy} after an onset other than \vi{qu}, where the older convention marks the first letter (\vi{hòa}, \vi{khỏe}) and the newer one the second (\vi{hoà}, \vi{khoẻ}): 69 of the 6,611 lower-case entries of the standard spelling list \citep{hunspell-vi}. Both are in everyday use; the older one is the majority in the Vietnamese test set we re-encode (51 of 53 affected tokens) and, as far as we can tell, in running text, the newer one the textbook convention (\placeholder{corpus count on two registers before Gate~1 fixes the majority and the direction of \hyp{4}}).

\paragraph{Three encodings of the same letters.}
A Vietnamese letter carries up to two diacritics, a quality mark (\vi{ă â ê ô ơ ư}) and a tone mark. Unicode has a precomposed code point for every combination, so the same text can be stored in NFC (one code point per letter), in NFD (base letter plus combining marks, with the tone mark \emph{before} the quality mark in \vi{ộ ậ ặ ệ}) or in a partially composed form inherited from Windows-1258 (quality mark precomposed, tone mark combining) that survives in legacy web text. All three are valid and occur in the wild. A tokenizer that normalizes to NFC erases the difference; one that passes its input through receives a different token sequence for the same word; some implementations \emph{corrupt} combining marks, deleting them. Neither Gemma tokenizer we audited normalizes (Section~\ref{sec:setup}).

\paragraph{Nói lái.}
\vi{Nói lái} is the Vietnamese spoonerism: two syllables exchange components so that a new, often humorous or hidden, phrase appears. Writing a syllable as $(O, R, T)$, a variant is a subset of $\{O, R, T\}$ exchanged between the two positions; the empty subset is the identity, the full one plain reversal, and the six others are the six types of the tradition \citep{le-ho-1990-thu-choi-chu,nguyen-van-hiep-noi-lai}: \var{1} swaps the rimes (\vi{mèo cái} $\rightarrow$ \vi{mài kéo}; \vi{trò chơi} `game' $\rightarrow$ \vi{trời cho} `heaven-sent'); \var{2} swaps onset and rime, so the syllables trade places and the tones stay (\vi{đầu tiên} `first' $\rightarrow$ \vi{tiền đâu} `where's the money?'); \var{3} swaps the tones (\vi{bí mật} `secret' $\rightarrow$ \vi{bị mất} `got lost'); \var{4} swaps rime and tone (\vi{bí mật} $\rightarrow$ \vi{bật mí} `reveal'); \var{5} swaps the onsets (\vi{giải pháp} `solution' $\rightarrow$ \vi{phải giáp} `must face') and \var{6} onsets and tones. A textbook illustration runs all six on \vi{thay đổi} `to change' (\vi{thôi đảy}, \vi{đôi thảy}, \vi{thảy đôi}, \vi{thổi đay}, \vi{đay thổi}, \vi{đảy thôi}), each of which our engine reproduces (\placeholder{confirm against \citet{le-ho-1990-thu-choi-chu}, not yet consulted directly}). Each variant is an involution, the eight form the group $(\mathbb{Z}/2)^3$, and reading an output in the other syllable order maps \var{1} to \var{6}, \var{2} to \var{3} and \var{4} to \var{5}; we therefore generate items for \var{1}--\var{4}, score \task{T1} in the named order and as an unordered pair, and keep \var{5} and \var{6} in the taxonomy, the \task{T2} gold and the \task{T3} twins (Section~\ref{sec:tasks}). When components coincide the variants collapse: \vi{thi đua} `to compete' $\rightarrow$ \vi{thua đi} `just lose' has equal tones, so \var{1} and \var{4} agree and \var{3} is the identity. Attested practice bends outputs toward real words: \vi{thầy giáo} `teacher' $\rightarrow$ \vi{tháo giày} `take off shoes' keeps the \vi{ay} of the real word where the rule gives \vi{giầy}, \vi{Vũ Như Cẩn} $\rightarrow$ \vi{vẫn như cũ} `still the same' is exact only under the Southern merger of \vi{hỏi} and \vi{ngã}, and three-syllable phrases swap a non-adjacent pair (\vi{chà đồ nhôm} `polish aluminium' $\rightarrow$ \vi{chôm đồ nhà} `steal from home'). The literature attaches regional labels to the variants inconsistently, so our codes are region-neutral and validators are region-tagged. Nói lái runs from folk sayings to the poetry of Hồ Xuân Hương, where it hides what a censor would not print; some attested examples have vulgar readings, which our release flags and keeps out of every human form and API prompt (see the Ethical Considerations section).\footnote{\placeholder{[NATIVE-CHECK] Every Vietnamese example in this paper is a row of the released attested seed table or an engine output on a dictionary pair; all forms and glosses await verification by the native validators.}}

% Section 3 -- The NoiLai benchmark. Budget: 1.50 pages, holds Table 1.
% Follows docs/DESIGN_DECISIONS.md sections 3-5, 10 and 11 (strict/lenient T1, six-variant
% T2 gold, spelling twins out of the T3 headline, filters, 4,200-item sample, canary header,
% attested exact/approx tags, validation and human-baseline protocol, release terms).
\section{The \noilai{} Benchmark}
\label{sec:benchmark}

\noilai{} is generated and scored by a rule engine that encodes the syllable structure and spelling rules of Section~\ref{sec:background}. We describe the engine, the three tasks, the sampling and splits, the attested subset, and the human validation. Table~\ref{tab:counts} gives the counts; every headline size is reported with its cluster count (items from base pairs).

% GENERATED by paper/gen_table1.py -- do not edit by hand.
% source manifest: /tmp/claude-0/-home-user-ntcf/8fecdaa4-9b4a-59bf-8d00-9f4658050a5a/scratchpad/paper_build/manifest.json
% generator commit fd27363bc105 (dirty tree), seed 20261004, generated 2026-09-30 15:10Z
% smoke build; regenerate: generator arguments differ from the plan's defaults or the build is smaller than the paper's ~10,000 items. Re-run this script on the frozen release.
\begin{table}[t]
\centering\small
\begin{tabular}{@{}llrrrrr@{}}
\toprule
Task & Variant & Dev & Test & Core & Total & Pool \\
\midrule
\multirow{5}{*}{\task{T1} transformation} & \var{1} & 27 & 123 & 20 & 150 & 328 \\
 & \var{2} & 32 & 118 & 20 & 150 & 394 \\
 & \var{3} & 32 & 118 & 20 & 150 & 296 \\
 & \var{4} & 26 & 124 & 20 & 150 & 423 \\
 & all & 117 & 483 & 80 & 600 & \\
\midrule
\multirow{5}{*}{\task{T2} decoding} & \var{1} & 12 & 68 & 20 & 80 & 197 \\
 & \var{2} & 24 & 56 & 20 & 80 & 232 \\
 & \var{3} & 16 & 64 & 20 & 80 & 172 \\
 & \var{4} & 18 & 62 & 20 & 80 & 263 \\
 & all & 70 & 250 & 80 & 320 & \\
\midrule
\multirow{5}{*}{\task{T3} validity} & \var{1} & 24 & 96 & 20 & 120 & -- \\
 & \var{2} & 28 & 92 & 20 & 120 & -- \\
 & \var{3} & 16 & 104 & 20 & 120 & -- \\
 & \var{4} & 20 & 100 & 20 & 120 & -- \\
 & all & 88 & 392 & 80 & 480 & \\
\midrule
Generated & & 275 & 1,125 & 240 & 1,400 & \\
Attested & & -- & \placeholder{200--400} & -- & \placeholder{200--400} & \\
\quad of which flagged vulgar & & & \placeholder{TODO} & & & \\
\midrule
Base pairs & & \multicolumn{5}{l}{500 (split by pair; dev fraction 0.2)} \\
Validation $\alpha$ & & \multicolumn{5}{l}{\placeholder{$\alpha$ = TODO, 95\% CI TODO; 1,000 items, 200 triple-judged}} \\
\bottomrule
\end{tabular}
\caption{\noilai{} item counts by task, variant and split. Core is the balanced subset of the test split that API-served models see (\task{T3}: yes/no pairs kept together). Pool is the number of legal, non-identity items available per cell before capping (\task{T1}, \task{T2}); the drop rate follows from it. Attested examples are folk and literary \vi{n\'{o}i l\'{a}i}; items flagged vulgar never reach an API or the human-baseline form. \placeholder{smoke build; regenerate: these counts come from a reduced build (1,400 items, seed 20261004); the frozen release replaces them.}}
\label{tab:counts}
\end{table}


\subsection{Parser, speller and legality}
\label{sec:engine}
The parser maps an orthographic syllable to a structure (onset, glide, nucleus, coda, tone) in canonical symbols, and the speller maps a structure back to standard orthography, applying every spelling rule of Appendix~\ref{app:rules}. Parsing is a spelling round-trip: a string is a well-formed syllable only if the standard spelling of its structure reproduces it, so \vi{céo} (for \vi{kéo}) is rejected in strict mode and recognized, in lenient mode, as a misspelling of a known structure. Two attested alternations are accepted on input and canonicalized: the \vi{i/y} variation after a consonant (\vi{lý} for \vi{lí}) and either tone-mark placement; the emitted \vi{i/y} form follows usage (\placeholder{convention confirmed with the validators}). Legality is decided against an attested inventory, not generated from phonotactic rules: the Hunspell vi\_VN list in both placement styles \citep{hunspell-vi}, extended with word-list syllables that parse strictly and reuse an already attested rime (\placeholder{261} additions; the \placeholder{16} distinct list entries the parser does not parse, 32 lines over the two placement files, and the \placeholder{3} level-toned stop-coda entries the inventory rejects are listed and explained in the data statement). A structure is legal when its onset plus rime occurs in the inventory with some tone and its tone respects the stop-coda constraint; the generator never emits a syllable that fails this test, and every gold answer is re-parsed before it is written. One analytic decision deserves mention: we treat the \vi{u} of \vi{qu} as the glide, as phonology does, whereas school grammar treats \vi{qu} as an onset, and the two analyses give different swaps for a \vi{qu}- syllable. Because no attested example settles it, \vi{qu}- syllables are excluded from generated inputs and kept as a separate stratum pending a 20-pair native check.

\subsection{Three tasks}
\label{sec:tasks}
A \emph{base pair} is two syllables $(A, B)$. Writing $A = (o_A, r_A, t_A)$ for onset, rime and tone, the four generated variants are the component permutations of Table~\ref{tab:variants} in Appendix~\ref{app:rules}.

\textbf{\task{T1} (transformation)} gives a two-syllable input and the name and step-by-step definition of a variant, and asks for the output. The gold is unique. The \emph{strict} score requires the two output syllables' structures in the named order; the \emph{lenient} score, reported as a secondary number, accepts the unordered pair, so that a traditional onset-swap answer to a \var{4} item is wrong strictly and right leniently, and both numbers are reported.

\textbf{\task{T2} (decoding)} gives a nói lái form and asks what meaningful phrase it hides. Since every variant undoes itself, naming the variant would collapse \task{T2} into \task{T1}; the prompt therefore withholds it, and the gold is the set of every lexical reading of the input under any of the six variants in either syllable order. The base pair is always one such reading (asserted at build time); inputs with several readings are kept, and the number of readings is a covariate. Lexicality for the gold is dictionary membership plus native validation of every core item's gold set; the validated set is the primary gold for the core, and a legal answer that is a valid variant reading outside the gold set is binned as plausible-non-gold and its rate reported. The variant the model implicitly applied is scored over the three unordered classes $\{\var{1},\var{6}\}$, $\{\var{2},\var{3}\}$, $\{\var{4},\var{5}\}$, since the finer label is not identifiable from the string.

\textbf{\task{T3} (validity)} pairs each correct \task{T1} output with an invalid twin that shares the input and differs in one component type: the output of another of the six variants when its string is distinct, a changed onset, a changed rime, a changed tone, or a misspelling that violates a \vi{c/k}, \vi{g/gh} or \vi{ng/ngh} rule while keeping the structure. The model judges whether the candidate is the correct output of the named variant. Yes and no items come in pairs with the same input, so the set is balanced by construction and open models can also be scored by forced choice over the pair. Because spelling twins test orthographic convention rather than sub-syllabic access, the headline \task{T3} number excludes them and reports them as their own column; generated answers are also reported as balanced accuracy, so that a bias toward one answer cannot inflate the headline (\placeholder{observed yes-rate per model}). Placement variants are never twins, because both placements are valid.

\subsection{Sampling, strata and splits}
\label{sec:sampling}
Lexical base pairs are the two-syllable entries of a 74k-word list whose syllables are both attested (\placeholder{47,535} candidates; word-list typos and loanwords are dropped by the attestation test); pseudo-pairs are quota-sampled from attested syllables with a fixed seed so that their joint distribution over zero onset, glide, stop coda, spelling trigger and tone-pair class matches that of the lexical base pairs, so that pseudo and lexical items differ only in lexicality; the achieved shares, and any quota that could not be filled, are recorded in the build manifest. Filters run in a fixed order and their drop counts are recorded per cell: legality of every output syllable; the spelling round trip; outputs equal to the input or to its plain reversal (equal tones make \var{3} the identity and \var{2} a plain reversal; equal rimes make \var{1} the identity); duplicate output strings for one base pair, merged with the set of variant labels that produce them; a blocklist screen for vulgar readings, applied per syllable and per syllable pair in either order to outputs, \task{T2} readings and \task{T3} candidates, and to an input only when the input pair itself is a taboo phrase, since a per-syllable input screen would remove ordinary vocabulary and skew the strata; overlap with the attested set or the demonstration pairs; and the \vi{qu}-, \vi{p}-onset, marginal-rime and loanword exclusions. Each item carries covariates that enter every regression: lexicality of input and output, syllable frequencies in a reference corpus (the word-list count as a secondary column), presence of a glide, a zero onset or a stop coda, the spelling triggers applied (e.g.\ \vi{c}$>$\vi{k}), the tone pair and whether the tones, onsets or rimes coincide, and whether a syllable is affected by the placement convention. \task{T1} cells are balanced between lexical and pseudo inputs; \task{T2} is built from lexical pairs only. The pool is split \emph{by base pair}, so that every item derived from one underlying pair falls in one split, into a public development set (20\%) and a gated test set. Open models are run on a stratified \emph{main sample} of 4,200 test items (350 per task $\times$ variant cell, the size that resolves a 10-point cell contrast at 80\% power under the design effect the pilot measures; \placeholder{realized design effect from the pilot}) that contains the balanced \emph{core} of 125 items per cell (\task{T3}: 62 yes/no pairs), the only \noilai{} data an API-served model sees; a separately built placement-enriched set of pseudo-pair items, disjoint from the released splits by base pair and excluded from the main sample, every one of which changes under C2, serves the C2 arm. Every test file opens with a header record carrying the BIG-bench canary sentence and every test item carries the canary GUID and do-not-train flags \citep{2023-emnlp-main-308,srivastava-etal-2023-beyond}; the GUID is drawn outside the seeded generator, and the paper prints only its SHA-256 (Appendix~\ref{app:datastatement}) while the GUID itself lives in the gated README, the dataset card and the file header, so that a hit detects benchmark rather than PDF contamination. The build seed is not published: development and test items are drawn and split from one seeded stream, so a public seed would regenerate the gated test split; the released generator mints new instances, and the gated file, identified by its content hash, is the reference instance. A sealed split with a different seed and canary is regenerated at release, never sent to any API and never released; only its hash and counts are published. Contamination is a weak threat to a generator that can mint a fresh test set, but it is a real one for the attested subset, which we test with a guided-instruction probe \citep{2308.08493}.

\subsection{Attested examples}
\label{sec:attested}
Folk and literary nói lái are collected from published collections into a seed table with input, output, the engine-derived set of variant labels (\var{1}--\var{6}) that reproduce it, the swapped positions for three-syllable phrases, glosses, a confidence grade, a vulgarity flag and the source, and the engine's own output for the same input and variant is recorded next to the attested form. Not every attested example is a pure rule output: some reach a real word only under a regional merger (\vi{vẫn như cũ} for the rule's \vi{vẫn như củ}: the Southern merger of \vi{hỏi} and \vi{ngã}), some bend a vowel or tone for meaning (\vi{độc hại} for the rule's \vi{đọc hại}: \vi{o}$\rightarrow$\vi{ô}), and some swap a non-adjacent pair in a three-syllable phrase. Rows carry a three-valued tag, \emph{exact}, \emph{approximate by a named regional merger} or \emph{approximate by a named substitution}, and a test requires every approximate row to be reachable from the rule output by exactly the named change; every row enters \task{T2} with the attested gold, but only exact two-syllable rows enter \task{T1}-style scoring and the memorization analysis (\hyp{6}), whose generated comparison items are pure rule outputs matched on variant, task, lexicality, spelling trigger and tone class. Attested inputs and outputs are reserved out of the generated pool. The seed holds \nAttested{} rows, \nAttestedExact{} reproduced exactly by the engine and \nAttestedApprox{} approximate, \nAttestedVerified{} of them native-verified (counts generated from the release file, never typed); the target is 200--400, every row verified by a native validator, and rows flagged vulgar never reach an API or a human form. Rows come from at least three collections and each is cited; a candidate found on the web enters only after the validators accept its reading and an author has checked the citation on the source page, and \hyp{6} is tested only if at least 100 native-verified exact two-syllable rows result, otherwise it is reported descriptively (\placeholder{verified exact rows at the data freeze}).

\subsection{Validation and human baseline}
\label{sec:validation}
Two or three adult native speakers (\placeholder{number and regions of the validators}), one per region where possible, validate the generator before any test-split run, after a keyed 16-item calibration round built from pairs outside the release. The core of the packet is a stratified probability sample of 30 items per task $\times$ variant cell (strata: cell $\times$ lexical or pseudo source, proportional allocation, weights recorded), mixed with four planted controls per cell, each a deliberately wrong candidate or verdict built by the engine from items outside the sample; the 408 rows are each judged by two validators, and 60 by all three. For every row the validator answers whether the output is the correct nói lái of the input under the named variant (for \task{T3}, whether the system's verdict is right), whether it is spelled acceptably, whether the input and the output are real Vietnamese words or phrases, and whether either has an offensive reading, and names the regional merger, if any, on which the judgment depends. Every validator also judges every attested row, a 20-pair \vi{qu}- check, 20 ``which nói lái do you produce?'' pairs that region-tag the variants, 14 spelling-convention questions and a share of the core \task{T2} gold sets (50 judged by all). Agreement is Krippendorff's $\alpha$ \citep{krippendorff-2011-computing,J08-4004}; $\alpha$ on correctness is computed over the sheet including its planted control rows, which supply the true negatives that $\alpha$ needs when nearly every item is correct, and the probability-sample-only $\alpha$ and AC1 are reported beside it, with raw agreement and the label marginals; MASI and Jaccard distances serve the set-valued \task{T2} gold. A disagreement goes to the third validator, who judges blind; a remaining split is decided by the authors against the rule tables with a logged reason, and the authors never override validators who agree. A confirmed generator bug triggers a fix and a full regeneration before any model run. Generator precision is a stratum-weighted estimate, pooled in the text and per cell in Appendix~\ref{app:validation}, where 30 items per cell give wide intervals; each validator's catch rate on the planted controls is reported. About twenty adult native speakers who are not validators each complete a 30-item form (30--45 minutes): six anchor items seen by everyone and 24 items assigned so that every other item is seen by exactly two people, 246 items in all drawn from the open-model main sample, so that every model is scored on exactly these items. Each row shows the exact p0 prompt the models receive, with its three demonstrations; one instruction-check row and a 10-item block of attested decoding items complete the form. The comparator is \emph{mean-human} accuracy (a model is one rater; ``any human correct'' is a ceiling reported separately), with an interval from a two-way bootstrap over persons and items; at $\pm$4--5 points it is a reference band, and no cell-level human comparison is claimed. Volunteers are unpaid adults who consent in writing and can withdraw; only coarse demographics are collected (the Ethical Considerations section and Appendix~\ref{app:validation}). \placeholder{Validation results: $\alpha$ (sheet and sample-only), raw agreement and AC1 per judgment with CIs, control catch rates, adjudication counts, pooled generator precision, number of items changed.}

\subsection{Release}
The generator, parser, re-encoder and scorer are released under Apache-2.0. The development split ships under CC~BY~4.0; the test split is gated and encrypted, with its canary, under CC~BY-NC-ND~4.0 (\placeholder{or GPLv2 with attribution if the word-list authors' permission has not arrived by the data freeze; DESIGN\_DECISIONS 4.1}). The GPL spelling and word lists are consulted at build time; the lexical items reproduce \placeholder{$N$} of their two-syllable entries, and the release says under which terms. Every build writes a manifest with counts and drop counts per cell, resource hashes, the content hash and the generator commit; the build seed is withheld (Section~\ref{sec:sampling}).

% Section 5 -- Experimental design: panel, re-encoding arms, metrics, pre-registered analysis. Budget: 1.20 pages.
% Setup budget line in main.tex. Taken from docs/DESIGN_DECISIONS.md sections 5-8 and docs/PREREGISTRATION.md
% sections 2-9, never redesigned (paper/claims.md section C). The model table and the full statistics are in
% Appendix app:analysis and app:compute.
\section{Experimental Design}
\label{sec:setup}

\paragraph{Panel.}
The planned panel of \placeholder{17} self-hosted and \placeholder{4} API-served models maximizes distinct tokenizers (Table~\ref{tab:models}): Gemma~3 1B/4B/12B \citep{gemma-team-2025-gemma3}, Qwen~3.5 and Gemma~4 size series, Southeast Asian fine-tunes, the Vietnamese-first PhoGPT-4B-Chat \citep{2311.02945} and Vistral-7B, Llama-3.1-8B, two 27B references and four hosted models without log-probabilities. The panel is frozen after smoke tests with pinned revisions; between-tokenizer claims need at least eight open models from four tokenizer families. No system prompt is used, thinking is off, and the chat template is tokenized once in the runner so that the hashed sequence is what every backend scores.

\paragraph{Prompts.}
Vietnamese instructions define the syllable structure, the variant and the spelling rules, with three fixed demonstrations (Appendix~\ref{app:prompts}) and three paraphrases, since formatting alone moves accuracy \citep{2310.11324}. Decoding is greedy (64 tokens); an output without an extractable answer counts as wrong. Giving each syllable already decomposed into onset, rime and tone is a main condition, the only one that separates rule knowledge from sub-token access; other ablations are in Appendix~\ref{app:analysis}.

\paragraph{Re-encoding arms.}
The baseline is NFC in the older placement convention. C1 rewrites the text into NFD; C1$'$ into the partially composed form (\texttt{pc}), which adds fewer tokens, so that the three encodings are ordered in token count; C2 moves the tone mark to the newer convention, changing exactly the 69 affected syllable types; C3a/C3b strip tone marks or all diacritics, as information-destroying contrasts. The meaning-preserving arms are applied to the whole prompt, so that the item is not the one oddly encoded span; applying them to the item text alone is a secondary condition (500 Gemma~3 1B items). A pre-flight \emph{census} classifies each tokenizer and engine path as passing NFD through, normalizing it (C1 is then zero by construction) or corrupting it (the arm is refused). Gemma~3 passes through (98.7\% of 1,000 census strings get new ids).

\paragraph{Metrics.}
The primary metric is strict \task{T1} accuracy after canonicalizing Unicode form, placement, case and \vi{lí/lý} (misspellings and homophones are not repaired), averaged over paraphrases and pooled over \var{1}--\var{4}. Every output receives one of eleven error classes from the same parse. \task{T3} is scored as accuracy and balanced accuracy; for open models we add, in separate columns never compared with the generated answer, a yes/no forced choice (\texttt{t3\_yesno\_lp}) and a BLiMP-style pair score under the shared \task{T1} prompt (\texttt{t3\_pair\_lp}), raw and length-normalized per token, since twins differ in token count \citep{2023-emnlp-main-306}. No metric uses an LLM judge \citep{2025-findings-emnlp-587}; baselines are copy, reversal, a wrong-variant rule, random legal output and the mean human. Each tokenizer's profile is joined to every item: tokens per syllable, single-token share and \emph{boundary alignment}, the share of token boundaries inside a split syllable that fall on a sub-syllabic seam (adapting \citealp{2026-acl-long-634}). For Gemma~3: 1.78 tokens per syllable, 29.5\% single-token, alignment 0.92 among split syllables (0.94 over all syllables, a single-token syllable counting as 1.0) under NFC, against 2.83, 7.1\%, 0.21 (0.27) and a tone mark isolated in 41.6\% of syllables under NFD.

\paragraph{Pre-registered analysis.}
The base pair is the cluster: every interval is a base-pair cluster bootstrap (2,000 replicates) \citep{2411.00640}, every comparison is paired by item inside each replicate, and the primary inference for a paired contrast is that bootstrap or a paired $t$ test on base-pair means; McNemar's mid-$p$ \citep{mcnemar-1947-note,P18-1128} is secondary and labelled as ignoring clustering. Holm's correction \citep{holm-1979-simple,Q17-1033} runs within declared families, for the re-encoding table the fourteen cells of each model's row (four arms $\times$ \task{T1}, \task{T3}, XCOPA, plus placement $\times$ \task{T1}, \task{T3}) minus undefined ones. A claim about models in general rests on a hierarchical estimate with models nested in tokenizer families (REML with the Hartung--Knapp adjustment), reported with the number of families. For \hyp{1}, item correctness is regressed on whether a syllable is split, tokens per syllable centred within model and alignment among split syllables, with lexicality, corpus frequency, variant and task, model fixed effects and model $\times$ fragmentation interactions; \hyp{1} is decided by a nested likelihood-ratio test of adding split and alignment. For the arms the estimand is the within-item effect $\tau_m(a)$ (Section~\ref{sec:counterfactuals}). Table~\ref{tab:hypotheses} lists all hypotheses. 350 items per cell resolve a 10-point paired contrast at 80\% power, the 125-item API cells 11--17 points, so cell-level claims rest on open models; a 200-item development pilot fixes the design effect before any test-split run (\placeholder{pilot ICC and discordance}). Only \placeholder{16\%; page-level check of arXiv 2511.04703 pending} of 445 reviewed LLM benchmarks report uncertainty or tests \citep{2511.04703}; every table here does. Details: Appendix~\ref{app:analysis}.

% Section 5 -- Results: which models can do noi lai, and what predicts failure.
% Budget: 1.50 pages, holds Table 2 and Figures 2-3.
% STATUS: hypotheses only, in the wording of docs/DESIGN_DECISIONS.md section 1. No model has
% been run. Every result is a \placeholder.
\section{Results: Who Can Do Nói Lái, and What Predicts Failure}
\label{sec:results}

\paragraph{Pre-registered hypotheses.}
\hyp{1} (\emph{alignment beyond count}): within models, alignment has a positive and tokens per syllable a negative coefficient, and adding the split indicator and alignment to a model with tokens per syllable improves fit, decided by the nested likelihood-ratio test of Section~\ref{sec:setup}; a model with fewer than 300 misaligned split syllables in the main sample contributes descriptively only, and if fewer than four tokenizer families clear that floor \hyp{1} is descriptive. \hyp{2} is a descriptive case study, not a tested hypothesis, of PhoGPT-4B-Chat's whole-syllable vocabulary. \hyp{6} (\emph{memorization}): exact attested items outscore generated items matched on variant, task, lexicality, spelling trigger and tone class, and a guided-instruction probe \citep{2308.08493} recovers attested but not matched items; it is tested only if at least 100 native-verified exact two-syllable rows from three collections exist at the data freeze.

\paragraph{Overall accuracy.}
\placeholder{Table~\ref{tab:main} reading: models above the human band and at the wrong-variant floor; variant ordering on non-degenerate items (expected \var{2} easiest, \var{3} hardest); \task{T3} against 50\%.}

\begin{table*}[t]
  \centering\small\setlength{\tabcolsep}{3.5pt}
  \begin{tabular}{@{}lrrrrrrr@{}}
    \toprule
    & \multicolumn{4}{c}{\task{T1} transformation, strict (by variant)} & \task{T2} & \multicolumn{2}{c}{\task{T3} (no spelling twins)} \\
    \cmidrule(lr){2-5}\cmidrule(lr){6-6}\cmidrule(lr){7-8}
    Model & \var{1} & \var{2} & \var{3} & \var{4} & decoding & balanced acc. & log-prob.\ (yes/no; pair) \\
    \midrule
    Mean human (\placeholder{$n$} raters, 246 items) & \placeholder{--} & \placeholder{--} & \placeholder{--} & \placeholder{--} & \placeholder{--} & \placeholder{--} & -- \\
    Copy / reversal & 0.0 & 0.0 & 0.0 & 0.0 & 0.0 & 50.0 & 50.0 \\
    Wrong-variant rule & \placeholder{--} & \placeholder{--} & \placeholder{--} & \placeholder{--} & \placeholder{--} & \placeholder{--} & -- \\
    \midrule
    \placeholder{per model} & & & & & & & \\
    \bottomrule
  \end{tabular}
  \caption{Accuracy (\%) with 95\% base-pair bootstrap intervals, mean over three paraphrases (range, lenient \task{T1} and spelling twins in Appendix~\ref{app:extra}). Human row: mean rater on the 246 double-judged items, on which every model is also scored. Log-prob.: open models only, never compared with the generated answer. \placeholder{TODO: filled by \texttt{noilai.stats.tables} from the scored runs.}}
  \label{tab:main}
\end{table*}

\paragraph{What predicts failure (E2).}
\placeholder{Figure~\ref{fig:alignment} reading; verdict on \hyp{1} (coefficients, likelihood-ratio test, floor counts, GEE agreement, share of per-model slopes with the predicted sign); the tokenizer-fixed SEA-LION pair (suggests, never isolates); \hyp{2}.}

\begin{figure}[t]
  \centering
  \fbox{\parbox{0.95\columnwidth}{\centering\vspace{1.8cm}\placeholder{Figure 2: accuracy vs.\ boundary alignment and tokens per syllable; coefficients with bootstrap intervals}\vspace{1.8cm}}}
  \caption{\placeholder{Item-level accuracy as a function of boundary alignment among split syllables (left) and tokens per syllable (right), pooled over models with 95\% intervals; coefficients from the fixed-effect fit with base-pair cluster-bootstrap intervals (inset).}}
  \label{fig:alignment}
\end{figure}

\paragraph{Explicit input.}
\placeholder{Decomposed vs.\ raw input per model: near ceiling places the bottleneck on access, low on the rule.}

\paragraph{Error taxonomy.}
\placeholder{Figure~\ref{fig:errors} reading against the registered expectations: spelling errors on trigger items, copy and reversal in the smallest models.}

\begin{figure}[t]
  \centering
  \fbox{\parbox{0.95\columnwidth}{\centering\vspace{1.8cm}\placeholder{Figure 3: error breakdown by model group}\vspace{1.8cm}}}
  \caption{\placeholder{Error classes as shares of all \task{T1} outputs, by model group and variant.}}
  \label{fig:errors}
\end{figure}

\paragraph{Ablations and memorization.}
\placeholder{Ablations; \hyp{6} with balance and the four discriminators.}

% Section 6 -- Orthographic counterfactuals. Budget: 1.00 page, holds Table 3 and Figure 4.
% Follows docs/DESIGN_DECISIONS.md sections 6 and 8.6 (per-model ATE, three-arm contrast,
% tone-isolation DiD, placement direction, no mediation language).
% STATUS: hypotheses only. No model has been run. Every result is a \placeholder.
\section{Orthographic Counterfactuals}
\label{sec:counterfactuals}

The arms change only the code points, hence the token sequence, so a difference between arms is caused by tokenization and by nothing else about the item. For arm $a$ on model $m$ the estimand is $\tau_m(a) = \mathbb{E}_i[Y_i(a) - Y_i(\mathrm{base})]$, identified by scoring the same items under both arms with all else fixed. C1 and C1$'$ run on \noilai{} and on the 500 Vietnamese test items of XCOPA \citep{ponti-etal-2020-xcopa}, pure NFC in the older convention; C2 runs on the placement-arm set only, since it touches 43 XCOPA items.

\paragraph{Pre-registered hypotheses.}
\hyp{4} (\emph{the rarer convention costs}) leads, because both placement conventions are common, so distribution shift cannot explain a C2 effect: the convention rarer in a reference corpus (expected: the newer; fixed by a count before any run) lowers accuracy on affected items, claimed as a pooled estimate with the number of models whose Holm-adjusted interval excludes zero and equivalence tests for the rest. \hyp{3} (\emph{a crossover}): in pass-through tokenizers $\tau_m(\mathrm{NFD})$ is positive on \task{T1} \var{3}, where NFD makes the tone an explicit symbol, and negative on XCOPA, where it only lengthens the text. The default explanation it must beat is distribution shift: NFD is rare text, and under Gemma~3 it is longer for 91\% of words, carries 59\% more tokens per syllable and has alignment 0.27, so a uniform drop is reported as out-of-distribution sensitivity. The per-family table is primary, and with fewer than five informative families the per-family intervals decide; normalizing tokenizers are the negative control. \hyp{3b}: in Gemma~3 the NFD gain on \var{3} is larger where the tone mark becomes its own token than on matched items where it stays attached, and the same contrast on \var{1} is near zero. Holm runs over the fourteen cells of each model's row.

\paragraph{Effects.}
\placeholder{Table~\ref{tab:arms} reading; verdicts on \hyp{4}, \hyp{3} and \hyp{3b} from the pooled estimates (43 affected XCOPA items descriptive, Appendix~\ref{app:extra}).}

\begin{table*}[t]
  \centering\small
  \begin{tabular}{@{}llrrrrrrr@{}}
    \toprule
    Model & Census & $\Delta$tok & \multicolumn{2}{c}{C1: NFD} & C1$'$: PC & C2: placement & C3a: no tones & XCOPA C1 \\
    & & C1 & \task{T1} all & \task{T1} \var{3} & \task{T1} all & \task{T1} affected & \task{T1} all & \\
    \midrule
    \placeholder{per model} & \placeholder{verdict} & & & & & & & \\
    \bottomrule
  \end{tabular}
  \caption{Within-item effect of each re-encoding arm on accuracy (percentage points, arm minus baseline), paired over identical items, with 95\% clustered intervals; bold: significant after Holm within the model's row. Census: passes through, normalizes (zero by construction) or corrupts (arm refused). \placeholder{TODO: filled by \texttt{noilai.stats.tables.intervention\_table}.}}
  \label{tab:arms}
\end{table*}

\paragraph{Does the effect track token count or the tone symbol?}
Token count is a coarsening of the intervention, not a mediator, so we do not estimate a mediated share. Four pre-registered analyses (Appendix~\ref{app:analysis}) ask whether count can account for the effect: effect modification by the token-count change (associational), a zero-dose contrast, the ordered contrast baseline $<$ PC $<$ NFD, and encoding $\times$ character spacing. \placeholder{Reading of Figure~\ref{fig:deltatok} (Appendix~\ref{app:extra}).}


% Section 7 -- Where tone lives. Budget: 1.00 page, holds Figure 5.
% Follows docs/DESIGN_DECISIONS.md section 9 (two-condition decodability, positions, the
% legal V3 patching pair, two readouts, position groups; bi mat / bi mat is a V4 pair only).
% STATUS: hypotheses only. No model has been run. Every result is a \placeholder.
\section{Where Tone Lives}
\label{sec:tone}

Behavioural failure on \var{3} can arise because the model never represents the tone of a syllable it received in pieces, because it represents it but does not move it to the answer position, or because it moves it and cannot use it. The two Gemma~3 models that run in fp32 on a free GPU let us tell these apart.

\paragraph{Probes.}
Following \citet{2025-findings-emnlp-719}, we train layer-wise linear probes on the residual stream to decode a syllable's tone (and, secondarily, onset, coda, nucleus and rime) at four positions: the onset fragment of a split syllable, the syllable's last token, the bare tone-mark token that NFD sometimes produces, and the first token after the syllable, the one position where NFC and NFD are compared on equal terms. Stimuli are legal syllables balanced by tone, including legal non-word syllables so that tone is not confounded with lexical identity, in several carrier sentences with a tone-neutral word after the slot; train and test syllables are disjoint, with a second split held out by rime. Tone counts as \emph{linearly decodable} at a layer and position only when two conditions hold: \emph{selectivity}, accuracy at least $\delta_{\mathrm{sel}}$ above a control probe trained on random labels per syllable type \citep{D19-1275}, with $\delta_{\mathrm{sel}} = 0.15$ provisional and frozen from the pilot; and \emph{excess} over a structural baseline, the same probe fitted on one-hot token ids plus the coda class, because stop codas admit only two tones and, under NFC, token count itself predicts tone. Probes show that information is present, not that it is used \citep{2022-cl-1-7}; we report a layer-0 baseline curve, and, on syllables that are single tokens, the layer at which tone accuracy first rises to half its maximum gain over the embedding, the tone counterpart of the reconstruction the character-level work reports.

\paragraph{Patching and steering.}
Activation patching locates where transplanted states change the output \citep{meng-etal-2022-locating,heimersheim-nanda-2024-use}. A pair is two prompts identical token for token except one syllable whose tone differs, both syllables attested and of equal segmental make-up; the varying syllable is the \emph{second} input syllable, because under \var{3} the first output syllable takes its tone and the two gold answers then diverge at their first token (clean \vi{bí mà} $\rightarrow$ \vi{bì má}, corrupt \vi{bí mạ} $\rightarrow$ \vi{bị má}; readout \vi{bì} vs.\ \vi{bị}, single pieces; pseudo-phrases, \placeholder{[NATIVE-CHECK]}). Pairs are kept only if their tokenizations align, the readout pieces differ in the tone mark alone, and the model's clean answer is correct with a logit gap of at least one nat; a model that loses every pair at the last filter has no manipulation readout, which is itself a result. We cache the clean residual stream, run the corrupt input with layer $\ell$ overwritten at a position group, and read the logit difference between the two answers at the first diverging token, normalized by the clean--corrupt gap. Two readouts separate perception from manipulation: a \emph{perception} readout asks which tone the syllable carries and reads a digit; the \emph{manipulation} readout is the \var{3} prompt itself. Position groups are the varying syllable and each of its tokens separately, the fixed syllable, the instruction, the answer marker and the final position; patching everything from the varying syllable onward must recover fully. The $2 \times 2$ of readout by source position distinguishes intact perception, intact manipulation, a copy-type failure (the answer tracks the first input syllable), and an ignore-type failure, in which the perception readout recovers where the manipulation readout does not. A difference-in-means direction between two tone classes, the baseline of CausalGym \citep{2024-acl-long-785}, tests steering against a random direction of equal norm and an onset-class direction that must not flip tone. Gemma Scope~2 sparse autoencoders are optional and must beat the probe baseline to be reported \citep{kantamneni25a}.

\paragraph{Pre-registered hypothesis.}
\hyp{5} (\emph{perception without manipulation}): in a Gemma~3 model whose \var{3} accuracy is below 25\%, tone meets both decodability conditions at a non-local position, the token after the syllable, at some layer in the first half of the network; and patching shows the perception readout recovering where the manipulation readout does not. Such a dissociation places the failure in manipulation rather than perception. Under NFD, tone is expected to become decodable at an earlier layer than under NFC because the tone mark is its own token. A clean localization without the dissociation would be equally informative.

\paragraph{Results.}
\placeholder{Figure~\ref{fig:probes}: layer-wise selectivity and excess for tone at the four positions under NFC and NFD in Gemma~3 1B and 4B, with the layer-0 baseline; the patching heatmap (layer $\times$ position group) for both readouts with recovery. Text: the first decodable layer per model, position and encoding; the patching bands for perception and manipulation; steering flip rates against the two null directions; pair retention at every filter; verdict on \hyp{5}.}

\begin{figure}[t]
  \centering
  \fbox{\parbox{0.95\columnwidth}{\centering\vspace{2.5cm}\placeholder{Figure 5: probe selectivity and excess by layer (NFC vs.\ NFD, four positions) and patching heatmap for the two readouts, Gemma~3 1B and 4B}\vspace{2.5cm}}}
  \caption{\placeholder{Left: probe selectivity (accuracy minus control) and excess over the structural baseline for tone at each layer of Gemma~3 1B and 4B, under NFC and NFD, at the syllable's last token and at the token after it. Right: recovery of the clean logit difference when the residual stream at layer $\ell$ and a position group is patched from the clean run, for the perception and the manipulation readout.}}
  \label{fig:probes}
\end{figure}

% Section 8 -- Discussion and conclusion. Budget: 0.50 page.
% STATUS: framing written; every conclusion that depends on a result is a \placeholder.
\section{Discussion and Conclusion}
\label{sec:discussion}

\placeholder{Answers to \hyp{1}--\hyp{6} in three sentences; implications for diacritic-heavy scripts, character knowledge and evaluation.}

Three points hold whatever the numbers. A benchmark whose gold is computed can be regenerated after a leak and scored without an LLM judge, which keeps comparisons paired. The re-encodings are free interventions on any deployed model, and the census that precedes them documents whether a class of Vietnamese text reaches a model as intended. And the token-count analyses are descriptive: we have resisted calling it mediation, and the causal claim stays with the item-level contrast. Steering along a tone direction and sparse-autoencoder features are left to future work.


% Mandatory, outside the page count; may not introduce new results (ARR CFP).
% Limitations -- mandatory, outside the page count, may not introduce new results (ARR CFP:
% "should not introduce new methods, analysis, or results"; "methodological caveats that the
% reader should be aware of when interpreting the reported findings").
% From docs/DESIGN_DECISIONS.md 11.7, 10.2, 8.4, 3.5 and the founding plan (section 2.7);
% items that depend on results are marked and must be re-read against the final numbers.
% Unnumbered: no \label (it would resolve to the preceding numbered section); refer to it by name.
\section*{Limitations}

\paragraph{Generated items are not natural usage.}
\noilai{} items are produced by a rule engine from real and pseudo base pairs. A generated output is, by design, often a phrase that no Vietnamese speaker would say; the benchmark measures whether a model can execute a component permutation and spell the result, not whether it appreciates wordplay. We therefore call the generated items rule-generated component permutations and reserve the word nói lái for the attested set, which is the only natural-usage data, and which is small and, as we note below, contaminated. Human validation checks that the gold is correct, not that the items are typical, and the human baseline measures natives following linguistic instructions on such permutations, which is why it is scored three ways and never compared cell by cell.

\paragraph{The taxonomy's regional labels are inconsistent, and the spelling analysis is ours.}
The six-type taxonomy is documented in newspaper and popular sources and in one book we could not consult directly; the regional labels those sources attach to the variants contradict one another, so our variant codes carry none and validators are region-tagged instead. Our spelling rules are anchored in published grammars and verified against the two released spelling lists, but the analytic choices they encode (the \vi{qu} glide, the \vi{gi} contraction, the treatment of \vi{ă} before semivowel codas, the emitted \vi{i/y} form) are ours, and a native speaker may disagree with the swap a different analysis would predict. We excluded the one stratum (\vi{qu}-) where the disagreement produces different gold answers, which removes those syllables from the benchmark rather than settling the question.

\paragraph{One author and three validators shape the gold.}
The engine was written by one author, and three volunteers judge a thousand-item probability sample. Regional variation in tone (the Southern merger of \vi{hỏi} and \vi{ngã}) means that some attested forms are correct in one region and not another; our gold records both forms where we know of them and tags approximate rows, but coverage of regional variants is not systematic. Agreement is reported as Krippendorff's $\alpha$ with raw agreement and Gwet's AC1 beside it, and generator precision as an estimate with an interval per cell; neither number is a guarantee about the items no validator saw.

\paragraph{The attested subset is small and contaminated.}
Folk nói lái examples circulate widely on the Vietnamese web, and the models were very likely trained on many of them. We treat this as a quantity to be measured (\hyp{6}) rather than a bias to be corrected, and we test it with a guided-instruction probe; but the attested numbers cannot be compared with the generated numbers as if they measured the same skill, and salience and humour cannot be matched away. The subset is also far smaller than the generated set, its sources are few, and if fewer than 100 verified exact rows exist at the data freeze the memorization hypothesis is reported descriptively only.

\paragraph{Cross-tokenizer comparisons are observational.}
Models differ in data, size, training recipe and tokenizer at once. Our alignment audit relates tokenizer properties to accuracy across models, and the regression controls for the covariates we measured, but no regression across 21 models identifies a causal effect of the tokenizer; syllable frequency is measured from a reference corpus that is a proxy for training exposure, not the models' training data, and is nearly collinear with whole-token status, so residual confounding remains; and a tokenizer-fixed pair of models is a single case. The causal claims of this paper are confined to the item-level interventions of Section~\ref{sec:counterfactuals}, where meaning and information are fixed and only the code points change, and to the patching experiments of Section~\ref{sec:tone}. Even there, a re-encoding changes the token sequence in several ways at once (count, boundaries, the presence of an explicit tone symbol, and how far the text is from the training distribution); the decompositions of an intervention effect by token-count change are associational, we do not estimate a mediated share, distribution shift is the default explanation a crossover must beat, and the C3 arms do not preserve meaning and are contrasts, not counterfactuals. The placement arm has almost no lever on natural text (0.5\% of XCOPA's syllable tokens), so it is not tested on natural text (its 43 affected XCOPA items are reported descriptively only) and its estimate comes from the separately built C2-enriched \noilai{} set.

\paragraph{API models drift, and their data terms differ.}
The four hosted models are queried through free tiers that may change the served model without notice. A provider whose tier uses submitted content for training never receives test items: it runs either on a paid key with a data-use opt-out or on a development-derived set of the same shape, reported separately and never in the main table, and the core goes only to providers with no-training terms. Model identifiers and dates are pinned, raw responses saved, and at most the core is ever sent, never the sealed split; an API result cannot be reproduced exactly later. Hosted models also return no log-probabilities, so their \task{T3} scores rest on generated answers alone, and gpt-oss models cannot switch reasoning off entirely.

\paragraph{Precision and hardware.}
Self-hosted models run in fp16, fp32 or 4-bit quantization on free T4 GPUs and in bf16 on a TPU; quantized and fp16 inference differs from the bf16 the models were trained in. About 200 items per quantized model are re-run in bf16 and the drift reported in Appendix~\ref{app:extra}, but the main numbers are those of the configuration that fits the hardware.

\paragraph{The mechanistic study covers two small models.}
Probing and patching are run on Gemma~3 1B and 4B only, the models whose fp32 weights fit free GPUs and whose every layer Gemma Scope~2 covers. Whether the localization generalizes to larger models, to other families, or to the API models we cannot inspect is an open question. Probes show that information is present, not that it is used; patching shows that a location carries the information the answer depends on in the pairs we constructed, which differ in one tone inside one prompt template, and a model that fails every pair has no manipulation readout at all.

\paragraph{Power per cell and prompt sensitivity.}
The 125-item API cells resolve only differences of 11--17 accuracy points, so cell-level claims rest on the open models' 350-item cells, and the human baseline, at $\pm$4--5 points, is a reference band rather than a cell-level comparator; a three-rater design over more items was beyond our volunteer budget. Results are reported over three native-written paraphrases of every prompt, and formatting alone has been shown to move accuracies by tens of points; three paraphrases bound but do not eliminate that sensitivity. Prompted judgments can understate what a model's probabilities reveal, which is why open models are also scored by log-probability; for hosted models the understatement cannot be measured.

\paragraph{Scope.}
This paper studies one language and one script. Vietnamese is unusual in writing every syllable as a separate word with its tone marked by a diacritic; the NFC/NFD re-encoder and diacritic stripping transfer to other diacritic-heavy Latin scripts, whereas the parser, speller and placement arms encode Vietnamese phonotactics and spelling and do not; the conclusions about tone do not transfer to languages whose tones are unwritten.

\placeholder{Re-read every paragraph above against the final results; a limitation that a result has resolved is removed, and a limitation a result has revealed is added. No new results may be introduced here.}

% Ethical considerations (optional section, outside the page count) and the anonymized
% Reproducibility paragraph. Mirrors docs/CHECKLIST_DRAFT.md and docs/DESIGN_DECISIONS.md 11.
% Unnumbered: it carries NO \label (a \label after \section* resolves to the preceding numbered
% section) and is referred to by name ("the Ethical Considerations section").
% STATUS: protocol. Nothing described here has happened yet; the closing \placeholder forces a
% re-read against what was actually done.
\section*{Ethical Considerations}

\paragraph{Human participants.}
Three native-speaker validators and about twenty human-baseline respondents take part. All are adults who volunteer, read a bilingual (Vietnamese and English) consent form describing the task, its duration and their right to withdraw at any time, and are not paid; we follow the precedent of community-built resources whose contributors were volunteers \citep{2024-acl-long-620,2020-findings-emnlp-195} and say so plainly. Participants are recruited by personal invitation from the authors' community and stand in no dependency relation to the authors. We collect no names, e-mail addresses or free text that could identify a person with the judgments; validators are identified by a letter whose key is held separately and deleted after publication, and forms record only a region tag and coarse demographics (age band, region where the participant grew up, years lived in Vietnam, whether raised abroad) for the data statement \citep{Q18-1041}. No participant information is ever sent to an API. The study is not reviewed by an institutional review board, because the authors have access to none; we answer the corresponding checklist item ``No'' with this justification, and we judge the risk to participants minimal: the task is a language judgment on wordplay items, participation is anonymous, every item with a known offensive reading is excluded from the human forms, and validators are warned that some items may have vulgar readings. Participants under 18 are not recruited, so no parental-permission procedure is needed.

\paragraph{Offensive content and dual use.}
Nói lái is a vehicle for satire and for hiding taboo meanings, and a benchmark of it necessarily contains items whose output, or whose decoding, is vulgar. A blocklist of taboo syllables and syllable pairs, seeded by the authors and extended by the validators, is applied at generation to gold answers, \task{T2} readings and \task{T3} candidates, per syllable and per syllable pair in either order, because ordinary inputs can have obscene swaps, and to an input only when the input pair itself is a taboo phrase, since several blocklisted syllables are ordinary words in context; attested items are screened by hand. Flagged items are excluded from the public development split, from every prompt sent to a hosted model and from the human forms; they stay, flagged, in the gated test file so that scoring is complete, and the validators' offensive-reading judgments on the sample give the residual rate. A model that is fluent in nói lái could, in principle, slip taboo content past a keyword filter; the generator is a few hundred lines of rules any speaker could write, the same engine is the detector (decode both swaps and check the blocklist), and gated release is the mitigation.

\paragraph{Data and terms of use.}
The spelling list and the word list consulted at generation time are GPL-licensed; the release states which of their entries the lexical items reproduce and under which terms (\placeholder{written permission from their authors, or GPLv2 with attribution}). XCOPA is CC~BY~4.0 and is used under that license; its re-encoded variants are produced at run time and never stored. The development split is released under CC~BY~4.0 and the test split is gated, encrypted and released under CC~BY-NC-ND~4.0 with a canary string; the per-item score files are public, while raw model outputs, which quote test items, are released only inside the gated repository under the same terms. Three hosted providers' free tiers are used; all require users to be 18 or older. A provider whose tier uses submitted content for training never receives test items: it runs either on a paid key with a data-use opt-out or on a development-derived set of the same shape, reported separately and never in Table~\ref{tab:main}, and the core goes only to providers whose terms exclude training on inputs, recorded per provider in the run manifests. All API experiments run on the account of an eligible adult, which the run manifests record by role; at most the core items are ever sent, and the sealed split never is. We use the providers' services to evaluate their models, which their terms permit, and use no output to train a model.

\paragraph{Generative assistance.}
Items and gold answers in \noilai{} come from a rule engine, not from a language model, and no metric uses an LLM judge. Code in the released repository, its documentation and drafts of this manuscript are produced with an AI coding assistant under the authors' direction, then run, tested and reviewed by the authors, who are responsible for every line and every claim; the research questions, design, rule tables and analyses are the authors'. The use is disclosed in the Responsible NLP checklist and will be stated in the Acknowledgements of the final version, so that the disclosure does not compromise anonymity now. Every reference is checked by hand against its source before submission.

\paragraph{Reproducibility.}
The generator, parser, re-encoder, prompt templates, scoring code, statistics and the probing and patching code are released at \anonurl, together with the development split, the gated test split, the attested seed table, the tokenizer-audit results and the per-item score files (item ids, covariates and scores); raw model outputs quote gated test items and are released only inside the gated test repository. Every dataset build and every model run writes a manifest with item counts and drop counts per cell (the dataset's build seed is withheld, because one seeded stream draws and splits the development and test items; its content hash is its identity), resource hashes, checkpoint identifier and revision, tokenizer and chat-template hashes, serving engine and version, requested and resolved precision, quantization, engine flags, prompt-file hashes, canary check, normalization census, hardware, wall time and GPU-hours; the manifests are released with the outputs. Compute totals \placeholder{TODO} GPU-hours on Tesla T4s and \placeholder{TODO} TPU-hours on a TPU v5e-8, all on free tiers, plus \placeholder{TODO} API calls; the log is released. Randomness is seeded everywhere; a build is byte-identical for a given seed and resource set. The hypotheses, metrics, exclusion rules, stopping rules and analysis plan are pre-registered before any test-split run in a dated commit whose hash is given in Appendix~\ref{app:prereg}, and every deviation from it is logged with a date.

\placeholder{Re-read every paragraph above against what was actually done before submission: who took part and how they were recruited, which providers and accounts were used, the compute totals, the reference check and the pre-registration commit; every present-tense protocol sentence becomes a past-tense report or is corrected.}


% references.bib: verified entries (Anthology-id keys; docs/RELATED_WORK_VERIFICATION.md).
% references_PLACEHOLDER.bib: only the entries the verified file lacks, every one marked
% UNVERIFIED; each must be checked by hand and moved into references.bib before submission.
\bibliography{references,references_PLACEHOLDER}

\appendix
% Appendices. Reviewers need not read them; every claim must stand in the eight pages.
% Tables under tables/ are generated (gen_rule_tables.py, gen_prompt_appendix.py).
% Protocols follow docs/DESIGN_DECISIONS.md sections 10 and 11.

\section{Prompts}
\label{app:prompts}

Every prompt is rendered from a YAML template by the evaluation harness; the three paraphrases (p0, p1, p2) carry the same content in different wording and order, and the English-instruction ablation changes only the instruction language while the item and the answer marker \vi{đáp án:} stay Vietnamese. Demonstrations are built by the rule engine from three fixed base pairs (\vi{trung bình}, \vi{làm chủ}, \vi{vô hình}) whose syllables occur in no item; \task{T2} demonstrations are the same for every item because the prompt withholds the variant. No model receives a system prompt; the three demonstrations are inline in the single user turn. The examples below are the p0 templates with three demonstrations, exactly as rendered, for one example item per task; the sha256 of every prompt sent to a model is stored with its output. \placeholder{[NATIVE-CHECK] Every Vietnamese string awaits validation by a native speaker before the runs; the six-type overview and the region-neutral wording of DESIGN\_DECISIONS 3 are to be reflected in the templates.}

% GENERATED by paper/gen_prompt_appendix.py on 2026-09-30 17:05Z from prompts/noilai.yaml + prompts/demos.yaml; do not edit.
% prompt file sha256: demos.yaml=5fe710f99645, noilai.yaml=a390e6a44153, xcopa.yaml=de275fa2df4b, _all=002fb57ab3d6
\paragraph{Transformation (\task{T1}), variant \var{1}, paraphrase p0, three demonstrations.} Prompt id \texttt{vi-p0-s3-explained-raw}, sha256 \texttt{fc90f3c2182f}.
\begin{quote}\small
Nói lái là một cách chơi chữ trong tiếng Việt: hai âm tiết của một cụm từ trao đổi với nhau một số thành phần (phụ âm đầu, vần hoặc thanh điệu) để tạo thành một cụm từ mới. \\
Mỗi âm tiết tiếng Việt gồm ba phần: phụ âm đầu (có thể không có), vần (âm đệm nếu có, âm chính và âm cuối nếu có) và thanh điệu (một trong sáu thanh: ngang, huyền, sắc, hỏi, ngã, nặng). Ví dụ: \textquotedbl{}trung\textquotedbl{} gồm phụ âm đầu \textquotedbl{}tr\textquotedbl{}, vần \textquotedbl{}ung\textquotedbl{} và thanh ngang; \textquotedbl{}bình\textquotedbl{} gồm phụ âm đầu \textquotedbl{}b\textquotedbl{}, vần \textquotedbl{}inh\textquotedbl{} và thanh huyền.

Kiểu nói lái cần áp dụng: đổi vần, giữ phụ âm đầu và thanh điệu. \\
Cách làm từng bước: \\
1. Tách mỗi âm tiết thành phụ âm đầu, vần và thanh điệu. \\
2. Đổi vần của âm tiết thứ nhất cho vần của âm tiết thứ hai. \\
3. Giữ nguyên phụ âm đầu và thanh điệu của mỗi âm tiết ở vị trí cũ. \\
4. Ghép lại thành hai âm tiết và viết đúng chính tả. \\
Khi ghép lại, hãy viết đúng chính tả tiếng Việt: âm /k/ viết là \textquotedbl{}k\textquotedbl{} trước e, ê, i, viết là \textquotedbl{}q\textquotedbl{} trước âm đệm u và viết là \textquotedbl{}c\textquotedbl{} trong các trường hợp còn lại (ví dụ \textquotedbl{}c\textquotedbl{} ghép với vần \textquotedbl{}eo\textquotedbl{} phải viết là \textquotedbl{}keo\textquotedbl{}, không viết \textquotedbl{}ceo\textquotedbl{}); tương tự, \textquotedbl{}gh\textquotedbl{} và \textquotedbl{}ngh\textquotedbl{} đứng trước e, ê, i (ví dụ \textquotedbl{}nghe\textquotedbl{}, không viết \textquotedbl{}nge\textquotedbl{}); dấu thanh đặt đúng vị trí.

Ví dụ: \\
Cụm từ: trung bình \\
Đáp án: trinh bùng

Cụm từ: làm chủ \\
Đáp án: lù chảm

Cụm từ: vô hình \\
Đáp án: vinh hồ

Bây giờ hãy áp dụng kiểu nói lái trên cho cụm từ sau. \\
Cụm từ: mèo cái \\
Chỉ viết cụm từ kết quả trên một dòng riêng ở cuối, theo đúng dạng: \\
Đáp án: \textless{}cụm từ kết quả\textgreater{}
\end{quote}
\paragraph{Decoding (\task{T2}), paraphrase p0, three demonstrations (the variant is withheld).} Prompt id \texttt{vi-p0-s3-explained-raw}, sha256 \texttt{785246b4733d}.
\begin{quote}\small
Nói lái là một cách chơi chữ trong tiếng Việt: hai âm tiết của một cụm từ trao đổi với nhau một số thành phần (phụ âm đầu, vần hoặc thanh điệu) để tạo thành một cụm từ mới. Có bốn kiểu nói lái thường gặp: \\
- đổi vần, giữ phụ âm đầu và thanh điệu: Đổi vần của âm tiết thứ nhất cho vần của âm tiết thứ hai. Giữ nguyên phụ âm đầu và thanh điệu của mỗi âm tiết ở vị trí cũ. \\
- đổi cả phụ âm đầu và vần, giữ thanh điệu: Đổi chỗ phần phụ âm đầu và vần của hai âm tiết cho nhau (tức là đổi chỗ hai âm tiết). Giữ thanh điệu ở vị trí cũ: âm tiết đứng trước vẫn mang thanh của âm tiết đứng trước, âm tiết đứng sau vẫn mang thanh của âm tiết đứng sau. \\
- đổi thanh điệu, giữ phụ âm đầu và vần: Đổi thanh điệu của âm tiết thứ nhất cho thanh điệu của âm tiết thứ hai. Giữ nguyên phụ âm đầu và vần của mỗi âm tiết. \\
- đổi vần và thanh điệu, giữ phụ âm đầu: Đổi cả vần và thanh điệu của âm tiết thứ nhất cho vần và thanh điệu của âm tiết thứ hai. Giữ nguyên phụ âm đầu của mỗi âm tiết ở vị trí cũ. \\
Mỗi âm tiết tiếng Việt gồm ba phần: phụ âm đầu (có thể không có), vần (âm đệm nếu có, âm chính và âm cuối nếu có) và thanh điệu (một trong sáu thanh: ngang, huyền, sắc, hỏi, ngã, nặng). Ví dụ: \textquotedbl{}trung\textquotedbl{} gồm phụ âm đầu \textquotedbl{}tr\textquotedbl{}, vần \textquotedbl{}ung\textquotedbl{} và thanh ngang; \textquotedbl{}bình\textquotedbl{} gồm phụ âm đầu \textquotedbl{}b\textquotedbl{}, vần \textquotedbl{}inh\textquotedbl{} và thanh huyền.

Ví dụ: \\
Cụm từ: trinh bùng \\
Đáp án: trung bình

Cụm từ: chù lảm \\
Đáp án: làm chủ

Cụm từ: vồ hinh \\
Đáp án: vô hình

Cụm từ sau đây là một cách nói lái. Cụm từ gốc là gì? \\
Cụm từ: mài kéo \\
Chỉ viết cụm từ gốc trên một dòng riêng ở cuối, theo đúng dạng: \\
Đáp án: \textless{}cụm từ gốc\textgreater{}
\end{quote}
\paragraph{Validity (\task{T3}), variant \var{1}, paraphrase p0, a spelling twin as candidate.} Prompt id \texttt{vi-p0-s3-explained-raw}, sha256 \texttt{855d848b82af}.
\begin{quote}\small
Nói lái là một cách chơi chữ trong tiếng Việt: hai âm tiết của một cụm từ trao đổi với nhau một số thành phần (phụ âm đầu, vần hoặc thanh điệu) để tạo thành một cụm từ mới. \\
Mỗi âm tiết tiếng Việt gồm ba phần: phụ âm đầu (có thể không có), vần (âm đệm nếu có, âm chính và âm cuối nếu có) và thanh điệu (một trong sáu thanh: ngang, huyền, sắc, hỏi, ngã, nặng). Ví dụ: \textquotedbl{}trung\textquotedbl{} gồm phụ âm đầu \textquotedbl{}tr\textquotedbl{}, vần \textquotedbl{}ung\textquotedbl{} và thanh ngang; \textquotedbl{}bình\textquotedbl{} gồm phụ âm đầu \textquotedbl{}b\textquotedbl{}, vần \textquotedbl{}inh\textquotedbl{} và thanh huyền.

Kiểu nói lái đang xét: đổi vần, giữ phụ âm đầu và thanh điệu. \\
Cách làm từng bước: \\
1. Tách mỗi âm tiết thành phụ âm đầu, vần và thanh điệu. \\
2. Đổi vần của âm tiết thứ nhất cho vần của âm tiết thứ hai. \\
3. Giữ nguyên phụ âm đầu và thanh điệu của mỗi âm tiết ở vị trí cũ. \\
4. Ghép lại thành hai âm tiết và viết đúng chính tả. \\
Khi ghép lại, hãy viết đúng chính tả tiếng Việt: âm /k/ viết là \textquotedbl{}k\textquotedbl{} trước e, ê, i, viết là \textquotedbl{}q\textquotedbl{} trước âm đệm u và viết là \textquotedbl{}c\textquotedbl{} trong các trường hợp còn lại (ví dụ \textquotedbl{}c\textquotedbl{} ghép với vần \textquotedbl{}eo\textquotedbl{} phải viết là \textquotedbl{}keo\textquotedbl{}, không viết \textquotedbl{}ceo\textquotedbl{}); tương tự, \textquotedbl{}gh\textquotedbl{} và \textquotedbl{}ngh\textquotedbl{} đứng trước e, ê, i (ví dụ \textquotedbl{}nghe\textquotedbl{}, không viết \textquotedbl{}nge\textquotedbl{}); dấu thanh đặt đúng vị trí.

Ví dụ: \\
Cụm từ gốc: trung bình; cụm từ cần xét: trinh bùng \\
Đáp án: Có

Cụm từ gốc: làm chủ; cụm từ cần xét: chù lảm \\
Đáp án: Không

Cụm từ gốc: vô hình; cụm từ cần xét: vinh hồ \\
Đáp án: Có

Câu hỏi: Cụm từ \textquotedbl{}mài céo\textquotedbl{} có phải là kết quả đúng khi nói lái cụm từ \textquotedbl{}mèo cái\textquotedbl{} theo kiểu trên không? Kết quả đúng phải đúng cả phụ âm đầu, vần, thanh điệu và chính tả. \\
Trả lời \textquotedbl{}Có\textquotedbl{} hoặc \textquotedbl{}Không\textquotedbl{} trên một dòng riêng ở cuối, theo đúng dạng: \\
Đáp án: \textless{}Có hoặc Không\textgreater{}
\end{quote}


\section{Rule Tables}
\label{app:rules}

The tables in this appendix are generated from the parser's constants and the spelling lists, so that they cannot drift from the code: Table~\ref{tab:components} lists the canonical symbols, Table~\ref{tab:onsets} the surface spelling of every onset, Table~\ref{tab:spelling} the remaining spelling rules with engine-spelled examples, Table~\ref{tab:placement} the 69 syllables whose tone-mark letter differs between the two placement conventions, and Table~\ref{tab:variants} the six variants with engine-computed examples and the status of each (the four generated \task{T1} cells; \var{5}, the onset swap, and \var{6}, onsets and tones, are the order-reversed readings of \var{4} and \var{1} and enter the taxonomy, the \task{T2} gold and the \task{T3} twins).

% GENERATED by paper/gen_rule_tables.py (components) from noilai.vi.syllable constants and the Hunspell lists; generated 2026-09-30 15:54Z. Do not edit by hand.

\begin{table*}[t]\centering\small
\begin{tabular}{@{}lp{0.78\textwidth}@{}}
\toprule
Slot & Canonical symbols (as the parser writes them) \\
\midrule
Onset (24) & \texttt{$\varnothing$}, \texttt{b}, \texttt{c}, \texttt{ch}, \texttt{d}, \texttt{đ}, \texttt{g}, \texttt{gi}, \texttt{h}, \texttt{kh}, \texttt{l}, \texttt{m}, \texttt{n}, \texttt{ng}, \texttt{nh}, \texttt{p}, \texttt{ph}, \texttt{r}, \texttt{s}, \texttt{t}, \texttt{th}, \texttt{tr}, \texttt{v}, \texttt{x}. \texttt{c} stands for /k/ however spelled (c, k, q); \texttt{g} for the voiced velar fricative however spelled (g, gh); \texttt{ng} for the velar nasal however spelled (ng, ngh); \texttt{gi} is the onset that contracts a following i; \texttt{p} occurs only in loanwords. \\
Glide & present or absent (written \texttt{u} or \texttt{o}; the \texttt{u} of \texttt{qu} is the glide). \\
Nucleus (18) & \texttt{a}, \texttt{ă}, \texttt{â}, \texttt{e}, \texttt{ê}, \texttt{i}, \texttt{o}, \texttt{ô}, \texttt{ơ}, \texttt{u}, \texttt{ư}, \texttt{ia}, \texttt{iê}, \texttt{ua}, \texttt{uô}, \texttt{ưa}, \texttt{ươ}, \texttt{oo}. Open-only (never take a coda): \texttt{ia}, \texttt{ua}, \texttt{ưa}. Closed-only (always take a coda): \texttt{iê}, \texttt{uô}, \texttt{â}, \texttt{ă}, \texttt{ươ}. \\
Coda (11) & \texttt{$\varnothing$}, \texttt{m}, \texttt{n}, \texttt{ng}, \texttt{nh}, \texttt{p}, \texttt{t}, \texttt{c}, \texttt{ch}, \texttt{j}, \texttt{w}. \texttt{j} is the semivowel written i/y, \texttt{w} the semivowel written o/u; the stop codas \texttt{c}, \texttt{ch}, \texttt{p}, \texttt{t} take only the tones sắc and nặng. \\
Tone (6) & 0~=~ngang, 1~=~huyền, 2~=~sắc, 3~=~hỏi, 4~=~ngã, 5~=~nặng; index 0 (ngang) is the unmarked level tone, the other five carry a diacritic. \\
\bottomrule
\end{tabular}
\caption{The parser's canonical symbols (\texttt{noilai.vi.syllable}). A syllable is onset + rime + tone, the rime being glide$^?$ + nucleus + coda$^?$. Legality is decided against the attested inventory (Hunspell vi\_VN, both placement styles, extended with word-list syllables that reuse an attested rime).}
\label{tab:components}
\end{table*}

% GENERATED by paper/gen_rule_tables.py (onsets) from noilai.vi.syllable constants and the Hunspell lists; generated 2026-10-01 23:45Z. Do not edit by hand.

\begin{table*}[t]\centering\small
\begin{tabular}{@{}lll@{}}
\toprule
Onset & Written & Engine examples \\
\midrule
\texttt{$\varnothing$} & (nothing) & a, yên, ý \\
\texttt{b} & \texttt{b} & ba, bồng \\
\texttt{c} & \texttt{c, k, qu} & ca / ke / qua (c before a, k before e \^{e} i y, q before the glide) \\
\texttt{ch} & \texttt{ch} & cha, chồng \\
\texttt{d} & \texttt{d} & da, dang \\
\texttt{đ} & \texttt{đ} & đa, đồng \\
\texttt{g} & \texttt{g, gh} & ga / ghe (gh before e \^{e} i) \\
\texttt{gi} & \texttt{gi} & gia / gì / giêng (the rime's initial i is absorbed) \\
\texttt{h} & \texttt{h} & ha, hồng \\
\texttt{kh} & \texttt{kh} & kha, khang \\
\texttt{l} & \texttt{l} & la, lồng \\
\texttt{m} & \texttt{m} & ma, mồng \\
\texttt{n} & \texttt{n} & na, nồng \\
\texttt{ng} & \texttt{ng, ngh} & nga / nghe / ngoe (ngh before e \^{e} i; not with a glide) \\
\texttt{nh} & \texttt{nh} & nha, nhồng \\
\texttt{p} & \texttt{p} & pa, pin \\
\texttt{ph} & \texttt{ph} & pha, phồng \\
\texttt{r} & \texttt{r} & ra, rồng \\
\texttt{s} & \texttt{s} & sa, sồng \\
\texttt{t} & \texttt{t} & ta, tồng \\
\texttt{th} & \texttt{th} & tha, thang \\
\texttt{tr} & \texttt{tr} & tra, trồng \\
\texttt{v} & \texttt{v} & va, vồng \\
\texttt{x} & \texttt{x} & xa, xồng \\
\bottomrule
\end{tabular}
\caption{Surface spelling of the canonical onsets. The k/gh/ngh spellings are triggered by the first \emph{written} letter after the onset (e, \^{e}, i, y), so a glide switches them off (\texttt{ngoe}, \texttt{nguy}); every example is produced by \texttt{noilai.vi.syllable.spell}, and the examples of the plain rows are attested syllables of the release inventory.}
\label{tab:onsets}
\end{table*}

% GENERATED by paper/gen_rule_tables.py (spelling) from noilai.vi.syllable constants and the Hunspell lists; generated 2026-09-30 15:54Z. Do not edit by hand.

\begin{table*}[t]\centering\small
\begin{tabular}{@{}lp{0.42\textwidth}p{0.36\textwidth}@{}}
\toprule
Rule & Statement & Engine examples \\
\midrule
Glide letter & \texttt{u} after q and before \^{a} \^{e} \horn{o} \^{o} i ia i\^{e}; \texttt{o} before a \u{a} e & hoa, hoe, huê, tuân, qua \\
Nucleus i & written \texttt{y} after the glide and in the bare syllable; \texttt{i} elsewhere & quy, ý, lí, minh \\
Nucleus ia / i\^{e} & ya / y\^{e} after the glide; y\^{e} with a zero onset & quya, quyên, yêu, tiếng \\
Short \u{a} & written \texttt{a} before the semivowel codas & tay vs.\ tai; tau vs.\ tao \\
Coda /j/ & \texttt{y} after \u{a} \^{a}, \texttt{i} elsewhere & mây, mai, môi \\
Coda /w/ & \texttt{o} after a e, \texttt{u} elsewhere & chao, keo, mưu, chiêu \\
Stop codas & p t c ch take only sắc and nặng & mật, bất; *mât is rejected by the inventory \\
Open-only nuclei & ia ua \horn{u}a never take a coda & mía, mùa \\
Closed-only nuclei & i\^{e} u\^{o} \horn{u}\horn{o} \u{a} \^{a} always take a coda & muốn, mượn \\
Tone mark, new style & on the nucleus letter (second letter of i\^{e} y\^{e} u\^{o} \horn{u}\horn{o}, first of ia ua \horn{u}a) & tiếng, muốn, mía \\
Tone mark, old style & differs only on open oa oe uy after a non-q onset (mark on the first vowel letter) & hoà / hòa, khoẻ / khỏe, thuỷ / thủy \\
\bottomrule
\end{tabular}
\caption{Spelling rules the generator applies when it reassembles a syllable, with examples spelled by the engine. The same rules, run in reverse, are the parser's spelling round-trip: a syllable whose standard spelling does not reproduce the input is rejected in strict mode (so \vi{c\'{e}o} for \vi{k\'{e}o} scores as a spelling error, never as correct).}
\label{tab:spelling}
\end{table*}

% GENERATED by paper/gen_rule_tables.py (placement) from noilai.vi.syllable constants and the Hunspell lists; generated 2026-09-30 17:22Z. Do not edit by hand.

% 69 pairs (new style, old style)
\begin{table}[t]\centering\small
\begin{tabular}{@{}llllll@{}}
\toprule
New & Old & New & Old & New & Old \\
\midrule
choá & chóa & luỹ & lũy & toè & tòe \\
choé & chóe & ngoã & ngõa & toé & tóe \\
choẹ & chọe & ngoé & ngóe & toạ & tọa \\
chuỳ & chùy & nguỵ & ngụy & toả & tỏa \\
doá & dóa & nhoà & nhòa & toẻ & tỏe \\
doạ & dọa & nhoá & nhóa & toẽ & tõe \\
goá & góa & nhoè & nhòe & truỵ & trụy \\
hoà & hòa & nhoé & nhóe & tuý & túy \\
hoá & hóa & nhuỵ & nhụy & tuỳ & tùy \\
hoè & hòe & oà & òa & tuỵ & tụy \\
hoạ & họa & oé & óe & tuỷ & tủy \\
hoả & hỏa & oẹ & ọe & uý & úy \\
hoẹ & họe & oẻ & ỏe & uỷ & ủy \\
huý & húy & soà & sòa & xoà & xòa \\
huỷ & hủy & suý & súy & xoá & xóa \\
khoá & khóa & thoà & thòa & xoã & xõa \\
khoé & khóe & thoá & thóa & xoè & xòe \\
khoả & khỏa & thoả & thỏa & xoả & xỏa \\
khoẻ & khỏe & thuý & thúy & xoẹ & xọe \\
loà & lòa & thuỳ & thùy & xuý & xúy \\
loè & lòe & thuỵ & thụy & xuỳ & xùy \\
loé & lóe & thuỷ & thủy & đoá & đóa \\
luỵ & lụy & toà & tòa & đoạ & đọa \\
\bottomrule
\end{tabular}
\caption{The 69 syllables of the Hunspell vi\_VN list whose tone-mark letter differs between the new placement convention (mark on the nucleus letter) and the old one (mark on the first vowel letter of the open rimes oa, oe, uy). The set is the difference between the two released lists, and the engine's rule reproduces every pair; the C2 arm converts running text between the two conventions and changes exactly these syllables (and their capitalized forms).}
\label{tab:placement}
\end{table}

% GENERATED by paper/gen_rule_tables.py (variants) from noilai.vi.syllable constants and the Hunspell lists; generated 2026-10-01 23:45Z. Do not edit by hand.

\begin{table*}[t]\centering\small
\begin{tabular}{@{}lllll@{}}
\toprule
Variant & Swapped & Kept in place & Engine example & Status \\
\midrule
\var{1} & rime & onset, tone & \vi{mèo cái} $\rightarrow$ \vi{mài kéo} & generated cell \\
\var{2} & onset, rime & tone & \vi{đầu tiên} $\rightarrow$ \vi{tiền đâu} & generated cell \\
\var{3} & tone & onset, rime & \vi{bí mật} $\rightarrow$ \vi{bị mất} & generated cell \\
\var{4} & rime, tone & onset & \vi{bí mật} $\rightarrow$ \vi{bật mí} & generated cell \\
\var{5} & onset & rime, tone & \vi{giải pháp} $\rightarrow$ \vi{phải giáp} & taxonomy; \task{T2} gold; \task{T3} twins \\
\var{6} & onset, tone & rime & \vi{mèo cái} $\rightarrow$ \vi{kéo mài} & taxonomy; \task{T2} gold; \task{T3} twins \\
\bottomrule
\end{tabular}
\caption{The six \vi{nói lái} variants as component permutations of a syllable pair (\texttt{noilai.gen.variants}): the non-empty proper subsets of \{onset, rime, tone\} exchanged between the two positions. Each is an involution; reading an output in the other order maps \var{1} to \var{6}, \var{2} to \var{3} and \var{4} to \var{5}, so \var{5} and \var{6} add no new unordered pair and are not generated as \task{T1} cells. Examples \var{1}--\var{5} are attested folk forms that the engine reproduces exactly; the \var{6} example is the \var{1} pair read in the other order, and the seed's two attested \var{6} rows (\vi{thay đổi} $\rightarrow$ \vi{đảy thôi}, \vi{trái gió} $\rightarrow$ \vi{giái tró}) are reproduced exactly as well.}
\label{tab:variants}
\end{table*}


\section{Validation Protocol and Human Baseline}
\label{app:validation}

\paragraph{Validation.}
The packet has five parts, judged in order of priority. (A) Calibration: 16 items keyed by the rule engine on pairs that are not in the release (four plain correct \task{T1} items, five planted errors, and the cases the instructions single out: new-style placement, \vi{y} after a consonant, a vulgar output, a \vi{hỏi}/\vi{ngã} output, a \task{T2} item, a correct \task{T3} item and a \task{T3} spelling twin); a validator who matches the key on fewer than 80\% of the correctness answers is briefed again and takes a second calibration set; nobody is excluded on calibration. (B) Generated items: 30 per task $\times$ variant cell drawn as a stratified probability sample (strata: cell $\times$ lexical or pseudo source, proportional allocation by largest remainder; the population size, sample size and weight of every stratum are recorded; \task{T3} items are drawn singly, never both members of a twin pair) plus four planted controls per cell built from items outside the sample: a \task{T1} candidate replaced by another variant's output (never the gold in the other order, which is correct leniently) or by the gold with one tone changed, a \task{T2} candidate that is no reading of the input under any of the six variants in either order, and a \task{T3} item shown with the opposite verdict. The 408 rows are shuffled under opaque identifiers; 60 go to every validator and the rest to two validators in rotating pairs (all rows to both when there are two). Validators are told that some rows are deliberately wrong. (C) Every attested row, seed and web-sourced candidates alike, judged by every validator for validity, familiarity, spelling and offensiveness; a row counts as native-verified when at least two validators accept it and none rejects it, and it enters the attested set only with its citation checked on the source page. (D) Twenty pairs on which the validator writes the nói lái they produce first (onsets, rimes and tones all differ, so the answer identifies the variant), 20 \vi{qu}- pairs shown under both analyses, and 14 spelling-convention questions (tone-mark placement, \vi{i/y}, \vi{oao}, \vi{gi}). (E) Each core \task{T2} item's dictionary readings to accept or reject, with any missing reading added; 50 items are judged by every validator. For each Part B row the validator answers with \emph{yes}, \emph{no} or \emph{unsure}: (1) \emph{correct}: the output is the correct nói lái of the input under the named variant (for \task{T2}, the input is a nói lái of the candidate original under some variant and order; for \task{T3}, the system's verdict is right); (2) \emph{spelling}: the output is spelled acceptably, either tone-mark placement and the \vi{lí/lý} alternation being acceptable; (3) \emph{lexical}: the input, and separately the output, is a real Vietnamese word or phrase (the speaker's judgment, not dictionary membership); (4) \emph{offensive}: the input or the output has an offensive or vulgar reading, including one reached by decoding; and (5) \emph{dialect}: the regional merger, if any, on which the correctness answer depends. Each validator's region of upbringing is recorded. Agreement is Krippendorff's nominal $\alpha$ on yes/no with ``unsure'' as missing (as a third category in a sensitivity analysis), Gwet's AC1, raw pairwise agreement and the label marginals, each with a 1,000-replicate item bootstrap interval; $\alpha$ on correctness is computed over the Part B sheet including its planted control rows, and the probability-sample-only $\alpha$ and AC1 are reported beside it; Part E uses MASI and Jaccard distances. Adjudication: unanimous definite labels stand; two validators who disagree are resolved by the third, who judges blind, by strict majority; a remaining split, or a disagreement with only two validators, is decided by the authors against the rule tables with a logged reason; an item with fewer than two definite labels is unresolved, stays in the benchmark and is left out of the precision estimate; the authors never override validators who agree. One validator's offensive flag is enough to remove an item from every API prompt and the human forms. A validator who labels fewer than 75\% of the planted controls ``no'' is briefed again after the first weekly return, and the rate is reported; no validator's labels are excluded or down-weighted on it. A comment that points at a rule is checked against the rule tables, and a confirmed generator bug triggers a fix and a full regeneration before any model run, which is reported. Generator precision per cell is $\sum_h W_h \hat p_h$ over the cell's strata with population shares $W_h$, with an unweighted Wilson interval beside it; the pooled estimate weights cells by population. The instructions given to validators, in Vietnamese and English, the consent form and the adjudication log are released with the data. \placeholder{Validators: number, regions, hours; per-judgment agreement table; adjudication counts; generator precision per cell.}

\paragraph{Human baseline.}
About twenty adult native speakers who did not take part in the validation each complete one form of 30 items (30--45 minutes): six anchor items seen by everyone (one per generated variant, one \task{T2}, one \task{T3}) and 24 items assigned so that every non-anchor item is seen by exactly two people, 246 items in all, about twenty per cell, drawn from the open-model main sample so that every model is scored on exactly these items; items flagged vulgar are excluded. Each row shows the exact p0 user message the models receive for that item (instruction, the same three demonstrations, the item and the answer marker) with its prompt hash recorded, in randomized item order, with one instruction-check row and a closing 10-item block of \task{T2} items on attested, non-vulgar examples that serves as the human reference for the cultural task; respondents write the answer (\task{T1}, \task{T2}) or tick yes/no (\task{T3}). They are asked not to use a dictionary, a search engine or an AI assistant and to say whether they did, and whether they validated items for this study. Answers are scored by the same code as the models, strictly, leniently and tolerant of doublets and homophones. Item-level human accuracy is the mean of its two judgments; the comparator is mean-human accuracy, with ``any human correct'' as a ceiling reported separately and an interval from a two-way bootstrap over persons and items; results are reported by output lexicality and by region, and the agreement between the two raters of an item is reported as nominal $\alpha$ on the produced answer and on correctness. A form is excluded, and counted, when fewer than half of its items are answered, when the instruction check fails, or when the respondent reports using a tool or having validated. If more volunteers materialize, items receive a third rater before the item set widens.

\section{Data Statement}
\label{app:datastatement}

Following \citet{Q18-1041}. \emph{Curation rationale}: items are generated by a rule engine from two-syllable entries of a Vietnamese word list and from sampled attested syllables, to test whether models can manipulate the components of a syllable; attested examples are collected from published collections of folk and literary nói lái, each with its source. \emph{Language variety}: standard written Vietnamese (\texttt{vi-VN}), lowercase, NFC, stored in the majority placement convention, \vi{i/y} per the stated rule; the attested subset includes Southern forms. \emph{Speaker demographics}: not applicable to generated items; attested items are folk material of unknown authorship or literary quotations with their source cited; the attested curation reflects the author's region. \emph{Annotator demographics}: three adult native speakers and about twenty human-baseline respondents, described in aggregate by age band, region grown up in, years in Vietnam and whether raised abroad (\placeholder{reported at submission}). \emph{Speech situation}: written wordplay, out of context. \emph{Text characteristics}: two-syllable phrases (three in a few attested cases), no punctuation, with per-syllable structure, lexicality and frequency covariates. \emph{Provenance}: the spelling list and word list (GPLv2), XCOPA (CC~BY~4.0), per-row citations for attested items, the generator commit and content hash (the build seed is withheld), resource hashes; the SHA-256 of the canary GUID \placeholder{printed here at release}, the GUID itself living in the gated README, the dataset card and the file header; split sizes and vulgar-flag counts; the rule-versus-folk mismatches and out-of-design rows of the attested set; the list entries the parser rejects, listed and explained; known limitations (\vi{qu}, regional labels, the \vi{i/y} convention); and the full parse table of the syllable inventory as supplementary material. The long-form statement is in the repository.

\section{Pre-registration}
\label{app:prereg}

The pre-registration is two-stage. The generator, data build, sampling files, scoring, extraction, census, placement baseline and constants are committed first (commit \placeholder{stage-1 hash}); the hypotheses of Sections~\ref{sec:results}--\ref{sec:tone}, the metrics, exclusion rules, stopping rules, analysis plan and probing protocol are fixed before any test-split run (the pilot runs on development items only) and committed on \placeholder{date of the stage-2 commit} (commit \placeholder{stage-2 hash}); deviations are logged with a date in the repository. Table~\ref{tab:hypotheses} restates them with their directional predictions and the quantity that decides each.

\begin{table*}[t]
  \centering\small
  \begin{tabular}{@{}lp{0.44\textwidth}p{0.38\textwidth}@{}}
    \toprule
    & Prediction & Decided by \\
    \midrule
    \hyp{1} & Alignment carries information about \task{T1} accuracy beyond count: $\beta(\text{align}) > 0$, $\beta(\text{tokens}) < 0$, and adding the split indicator and alignment to a model with tokens per syllable improves fit. & Nested likelihood-ratio test / $\Delta$AIC (model fixed effects, base-pair cluster-bootstrap CIs for $\beta$); a model with fewer than 300 misaligned split syllables contributes descriptively only, and fewer than four families clearing the floor demotes \hyp{1} to descriptive. \\
    \hyp{2} & Descriptive case study, not a tested hypothesis: PhoGPT-4B-Chat's whole-syllable vocabulary is the panel's extreme on the share of whole-syllable tokens, an E2 covariate; its per-syllable NLL and \task{T1} accuracy are reported beside Qwen3.5-4B and Gemma~3 4B. & Reported with CIs; the whole-syllable-share coefficient in E2 is the population-level test; no inequality is pre-registered. \\
    \hyp{3} & In pass-through tokenizers, $\tau(\mathrm{NFD})$ on \task{T1} \var{3} $> 0$ and on XCOPA $< 0$; distribution shift is the default explanation the crossover must beat. & Two paired tests per model, Holm within the pair; both family-level pooled CIs on the predicted side, the per-family table primary; below five informative families the pooled CI is labelled unreliable and the per-family CIs decide; a uniform drop is reported as OOD sensitivity. \\
    \hyp{3b} & The NFD gain on \var{3} is larger where the tone mark is isolated than where it stays attached (DiD $> 0$); the \var{1} contrast $\approx 0$. & Matched items, cluster bootstrap by base pair (McNemar mid-$p$ as the secondary item-level statistic), Holm over the four contrasts. \\
    \hyp{4} & The rarer placement convention (expected: new style, fixed by a corpus count before any run) lowers accuracy on affected items in models generally. & Per-model paired test on affected items, Holm across the panel; $k$ of $N$ plus the pooled estimate; TOST with $\delta = 2$ points for nulls. \\
    \hyp{5} & In a Gemma~3 model with \var{3} accuracy $< 25\%$, tone is decodable (selectivity and excess) at the token after the syllable by layer $L/2$, and, on the pairs whose clean answer is correct, the perception readout recovers under patching where the manipulation readout does not. & Section~\ref{sec:tone} protocol; $\delta_{\mathrm{sel}} = 0.15$ provisional; syllable-clustered CIs; below 50 retained pairs, the perception readout alone and a localization study; no earlier-under-NFD claim where the structural baseline is near 100\%. \\
    \hyp{6} & Exact attested items outscore matched generated items (pooled); guided-instruction contamination $> 0$ on attested, $\approx 0$ on matched generated. & Matching on variant, task, lexicality, trigger and tone class; per-model tests Holm across the panel plus pooled; the contamination discriminators carry the inference; floor: at least 100 verified exact rows from at least three collections, else descriptive. \\
    \bottomrule
  \end{tabular}
  \caption{Pre-registered hypotheses, their directional predictions and the quantity that decides each; ``pooled'' is a hierarchical random-effects estimate with models nested in tokenizer families, reported with the number of families, never a count of significant models.}
  \label{tab:hypotheses}
\end{table*}

\section{Compute, Environment and Panel Details}
\label{app:compute}

\placeholder{Per-run table: model, checkpoint id and revision, engine and version, requested and resolved precision, quantization, hardware, wall time, GPU-hours; totals for checklist item C1. Tokenizer profile and census verdict of every panel model (tokens per syllable, single-token share, alignment among split syllables, onset$|$rime split, NFD tone isolation, passes through / normalizes / corrupts per engine). Package versions (C4). Models that failed smoke tests and why. The adult account holder's role for the API runs.}

\section{Additional Results}
\label{app:extra}

\placeholder{Accuracy range over paraphrases per model; lenient \task{T1}; the spelling-twin column of \task{T3}; bf16 drift on about 200 items per quantized model; the two-engine C1 agreement on Gemma~3 1B; the reasoning-cost curve (accuracy against reasoning tokens on 500 items); per-variant results restricted to non-degenerate items; per-stratum tables; the realized design effect and per-model ICC; the \vi{qu}- stratum if its convention is settled in time.}


\end{document}
